% !TeX program = lualatex
\documentclass[11pt,a4paper]{ltjsarticle}
\usepackage[margin=25mm,headheight=17pt]{geometry}
\usepackage[haranoaji]{luatexja-preset}
\usepackage{amsmath,amssymb,amsthm,mathtools,mathrsfs,booktabs,array,enumitem,longtable}
\usepackage[unicode,hidelinks]{hyperref}
\usepackage{fancyhdr}
\pagestyle{fancy}\fancyhf{}\fancyfoot[C]{\thepage}
\fancyhead[L]{\small\nouppercase{\leftmark}}
\renewcommand{\sectionmark}[1]{\markboth{#1}{}}
\setlength{\emergencystretch}{3em}\allowdisplaybreaks[2]
\setlength{\parskip}{.15em}
\newtheoremstyle{jp}{.6em}{.6em}{\normalfont}{}{\bfseries}{.}{.5em}{}
\theoremstyle{jp}
\newtheorem{theorem}{定理}[section]
\newtheorem{proposition}[theorem]{命題}
\newtheorem{lemma}[theorem]{補題}
\newtheorem{corollary}[theorem]{系}
\newtheorem{definition}[theorem]{定義}
\newtheorem{remark}[theorem]{注意}
\newtheorem{example}[theorem]{例}
\newtheorem{question}[theorem]{問題}
\renewcommand{\proofname}{証明}
\newcommand{\kk}{\mathbb C}\newcommand{\PP}{\mathbb P}
\newcommand{\OO}{\mathcal O}\newcommand{\A}{\mathcal A}
\newcommand{\B}{B_{\mathrm{VK}}}\newcommand{\Blo}{B_{p,q}}
\newcommand{\HH}{\operatorname{HH}}\newcommand{\Hom}{\operatorname{Hom}}
\newcommand{\End}{\operatorname{End}}\newcommand{\Sym}{\operatorname{Sym}}
\newcommand{\Perf}{\operatorname{Perf}}\newcommand{\Der}{\operatorname{Der}}
\newcommand{\Inn}{\operatorname{Inn}}\newcommand{\gr}{\operatorname{gr}}
\newcommand{\rk}{\operatorname{rank}}\newcommand{\Spec}{\operatorname{Spec}}
\newcommand{\id}{\operatorname{id}}\newcommand{\Ext}{\operatorname{Ext}}
\newcommand{\Tor}{\operatorname{Tor}}\newcommand{\Ann}{\operatorname{Ann}}
\newcommand{\coker}{\operatorname{coker}}\newcommand{\Supp}{\operatorname{Supp}}
\newcommand{\pd}{\operatorname{pd}}\newcommand{\depth}{\operatorname{depth}}
\newcommand{\Frac}{\operatorname{Frac}}\newcommand{\Aut}{\operatorname{Aut}}
\newcommand{\Out}{\operatorname{Out}}\newcommand{\GL}{\operatorname{GL}}
\newcommand{\PGL}{\operatorname{PGL}}\newcommand{\Hilb}{\operatorname{Hilb}}
\newcommand{\dashsim}{\mathrel{\overset{\mathrm{bir}}{\sim}}}
\setcounter{tocdepth}{1}

\hypersetup{pdftitle={六次Prym曲線の標準七次被覆と分岐曲線},pdfauthor={}}
\title{六次Prym曲線の標準七次被覆と分岐曲線}
\author{}\date{2026年10月29日\quad 第30版}
\begin{document}
\maketitle
\begin{abstract}
自己反対Fano72型関係代数の署名 $(3,3)$ の七次元商を、平面六次Prym曲線から復元する。
抽象的な曲線と非自明二捩れ線束の対は、平面表示、関係代数、二十一点配置を一意に決める。
係数Gram行列の核pencilが正則な開部分では、像は十九変数対称行列の逆行列に課す九つの零条件で記述される。
核pencilは標準的なPrym-canonical切断を定め、その十八点除数の五次イデアル成分は三次元である。
この三つの五次式は、平面から平面への内在的な次数七の有理写像を与える。
十八点を吹き上げた解消上で、一般の射は有限平坦であり、その幾何学的モノドロミーは全対称群 $S_7$ である。
分岐源は滑らかな幾何学的整曲線で種数三十七を持ち、標的の被約分岐曲線は幾何学的に整な二十四次曲線となる。
分岐源はその正規化であり、特異点の $\delta$ 不変量の総和は二百十六である。
正規化写像の微分零因子の長さは七十二、双対分岐曲線の次数は四十八となる。
証明では、収縮曲線の次数を三以下に制限し、同じ成分内の一次変形で三本の二次曲線を非有効化する。
さらに、滑らかな残余六次曲線を持つ特殊化を用い、三十節点の一次・二次平滑化を同じ成分内で検証する。
また、補助六次曲線の直接式、三つの通常二重点、および枠付き成分上のmonodromyの固定三点への $S_3$ 射影を記述する。
\end{abstract}
\begingroup\small\setlength{\parskip}{0pt}\tableofcontents\endgroup
\clearpage

\section{序論と主結果}

本稿の出発点は、同符号ブロックを含む全自己反対方程式である。
署名 $(3,3)$ では、三行七列の一次式行列 $P,B$ と非退化対称形式 $A$ が
平面六次曲線 $C=\{\det(PAP^t)=0\}$ を与える。
その対称余核から非自明二捩れ線束 $\eta$ が生じ、
二十一点配置のうち十八点は $|K_C\otimes\eta|$ の特定の除数となる。
本稿では、この特定の除数が抽象Prym対から内在的に決まり、
さらに平面の七次被覆を作ることを示す。

\begin{theorem}[復元、像の方程式と七次被覆]
\label{thm:sextic-main-v29}
標数零の代数閉体上で、指定反対合を保持した署名 $(3,3)$ の商 $\mathfrak Z_1^+$ を考える。
以下では、$C$ が滑らかで二十一点配置が被約となる非空開部分を取り、
第三項以降ではさらに $t$ の三つの成分が独立となる非空開部分へ制限し、
第四・五項はその中の一般点について述べる。
これらは同じ指定成分内の条件である。この範囲で次が成立する。
\begin{enumerate}
\item 写像 $\mathfrak Z_1^+\dashrightarrow\mathcal R_{10}$ は七次元の像への双有理写像である。
抽象対 $(C,\eta)$ から平面の直線束 $H$、関係代数、二十一点配置を復元できる。
二固有ブロックの順序は、全平面で階数三を持つブロックを選ぶことにより内在的に戻る。
\item 正則な係数Gram核pencilを標準化すると、枠付き像は十九変数対称行列 $A$ に対する
九本の余因子方程式で与えられる。
指定した整数基点を通る成分 $Y_0$ は十次元であり、残余三次元群の商が七次元のPrym像を与える。
\item 核pencilの標準直線から一次式ベクトル $t$ が定まり、十八非固定点 $D$ は
\[
 I_D=I_3[\,Q\mid t\,]=(\det Q,\operatorname{adj}(Q)t),\qquad Q=PAP^t,
\]
\[
 0\to S(-7)^3\to S(-5)^3\oplus S(-6)\to I_D\to0
\]
で記述される。
三つの一次式 $t$ が独立な非空開部分では、
三つの五次式 $g=\operatorname{adj}(Q)t$ の基点スキームは正確に $D$ である。
\item $X=\operatorname{Bl}_D\mathbb P^2$ 上でこの線形系は基点を持たず、
有限平坦な次数七の射 $\varphi:X\to\mathbb P^2$ を定める。
従って $L=5H-\sum_{i=1}^{18}E_i$ は豊富な線束である。
この七次元像の一般のPrym対について、その幾何学的モノドロミーは $S_7$ である。
この写像は、源と標的の射影座標変更を除き $(C,\eta)$ のみに依存する。
\item 分岐源 $R_\varphi$ は滑らかで幾何学的に整であり、
\[
 R_\varphi\sim12H-2\sum E_i=2\widetilde C,
 \qquad g(R_\varphi)=37
\]
を満たす。被約分岐曲線 $B_\varphi$ は幾何学的に整な二十四次曲線であり、
$\nu:R_\varphi\to B_\varphi$ は正規化である。従って
\[
 \sum_{p\in\operatorname{Sing}(B_\varphi)}\delta_p(B_\varphi)=216.
\]
さらに $d\nu$ の零因子の長さは七十二、双対曲線 $B_\varphi^*$ の次数は四十八である。
これらの数値には、分岐曲線の特異点を通常の節点と尖点に限定する仮定を用いない。
\end{enumerate}
\end{theorem}

十八単純基点を持つ五次netの自己交差が七になることは、標準的な交点計算である。
本稿の主眼は、そのnetを七次元の制約付きPrym像から標準的に取り出すこと、
その基点スキームを完全に決定すること、同じ成分内で有限平坦性と $S_7$ を証明すること、
および分岐源の滑らかさと分岐曲線の数値的幾何を決定することにある。
対称行列式表示と二捩れ線束の対応はBeauvilleの理論\cite{BeauvilleDeterminantal}を用い、
平面直線束の一意性は古典的Noether理論\cite{HartshorneNoether,CoppensKatoPlaneSystems}の範囲にある。
最小自由分解にはHilbert--Burch理論を用いる\cite{EisenbudSidman}。

ここで $\mathcal R_{10}$ は曲線と非自明二捩れ線束の対のモジュライであり、
Prym多様体だけからの復元を主張しない。
また標準七次写像の $S_7$ と、元の二十一点配置のモノドロミーは異なる被覆の群である。
十八非固定点の全モノドロミーや、反対合そのものの標識の除去は未決である。
署名 $(2,4)$ の四次Prym幾何・K3退化は別稿\cite{PaperV}に分け、
共通する自己反対方程式とGram曲線の基礎だけを本稿でも記す。

\section{全自己反対方程式と標識付き有理商}
\label{sec:prym-setup-v26}
基礎体は $\mathbb C$ とし、有限計算の入力は $\mathbb Q$ 上に取る。
$L(x)=\sum_{i=1}^3x_iL_i$ は六行七列の一次式行列である。
四頂点箙の矢の本数は $(3,6,3)$、左右の二次関係の次元はともに七であり、
右平面との同一視の後に右関係を $R_i=JL_i$ と書く。
\[
 J=\operatorname{diag}(I_p,-I_q),\qquad
 L=\begin{pmatrix}P\\B\end{pmatrix},\qquad A=A^t\in\operatorname{GL}_7.
\]
課す方程式は
\begin{equation}\label{eq:full-incidence-v26}
 L_iAL_j^tJ=JL_jAL_i^t\qquad(1\leq i,j\leq3)
\end{equation}
の\emph{全て}である。ブロックで書けば
\begin{align*}
 P_iAP_j^t&=P_jAP_i^t,& B_iAB_j^t&=B_jAB_i^t,\\
 P_iAB_j^t+P_jAB_i^t&=0 &&(i\leq j)
\end{align*}
となる。同符号の式を省いて混合直交性だけを使うと、係数 Gram 行列も逆構成も変わる。

対象の同値は、関係枠の $\operatorname{GL}_7$、二つの行ブロック内の
$\operatorname{GL}_p\times\operatorname{GL}_q$、共通の平面座標変更、
および pairing $A$ の非零スカラー倍で定める。
関係枠 $L\mapsto LM$ に対して $A\mapsto M^{-1}AM^{-t}$ とする。
有理商という語は、この作用の不変関数体を持つ幾何的な商の非空開部分を意味する。
普遍代数が粗モジュライ全体に降下するという主張は含めない。
$p=q=3$ では、まず $P$ に用いる固有空間を指定した商を使う。
二つのブロックの順序を忘れる操作は第\ref{subsec:unordered-blocks-v27}項で扱う。

\begin{proposition}[既存成分と開条件]
\label{prop:prym-baseline-v26}
別稿\cite{PaperII}で構成した固定成分 $Z_1,Z_2,Z_3$ のうち、
$p=3,2$ に対応する成分を、それぞれ $Z_1,Z_3$ と書く。
それらは、次の条件を同時に満たす非空開部分を持つ。
\begin{enumerate}
\item $P(x)$ は平面の全点で階数 $p$ を持ち、
$Z=V(I_6(L))$ は相異なる二十一点からなる。
\item 線形核から右関係を復元する写像の階数は所要の最大値であり、
pairing はスカラー倍を除いて一意である。
従って上の枠付き表示は、指定反対合を保つ関係代数の分類に対応する。
\item 点対合の固定点は $B$ の階数低下点であり、固定点数は
$f(p)=21-12p+2p^2$、非固定点数は $2pq$ である。
非固定点は $C$ 上にあり、固定点は $C$ の外にある。
\end{enumerate}
対応する標識付き有理商の次元は
\[
 \begin{array}{c|c|c|c|c}
 \text{成分}&(p,q)&\text{有理商}&\text{次元}&(\text{固定点},\text{非固定点})\\\hline
 Z_3&(2,4)&\mathfrak Z_3&8&(5,16)\\
 Z_1&(3,3)&\mathfrak Z_1^+&7&(3,18)
 \end{array}
\]
である。$A_{77}=1$ の全incidenceチャートは、それぞれ
百五十三変数と六十九式、百五十三変数と七十二式で定義され、
境界点 $B_-$ で滑らかである。
\end{proposition}

\begin{proof}[用いる入力と幾何学的内容]
関係成分の構成、商の次元および同じ成分での被約点配置の存在は
\cite[反対合の固定部分族・全署名の点配置の各節]{PaperII}を用いる。
付属の \nolinkurl{hensel_other_signatures.py} と二つの保存証明書は、
その有限入力を完全に含む。
Hensel 分枝は、滑らかな境界点で自由変数を指定し、単元 Jacobian を持つ残りの変数を一意に解いたものである。
本稿の新しい計算は、基点、pivot、自由変数を全て照合した同一分枝上で行う。
単に同じ素数を法として合同な別成分の点を選んでいるのではない。

計数の内容を確認しておく。混合直交性により、$B^t$ は階数 $q+1$ の束
\[
 \mathcal E=\ker\bigl(\mathcal O(1)^7\xrightarrow{PA}\mathcal O(2)^p\bigr)
\]
の $q$ 本の切断を与える。
その有限な階数低下スキームの長さは
\[
 c_2(\mathcal E)=[H^2]\frac{(1+H)^7}{(1+2H)^p}=21-12p+2p^2
\]
である。非固定点では $P$ と $B$ の行空間が非零に交わり、
直交性から $P$ の Gram 形式が退化する。
固定点が同時にその退化曲線上にあれば、$I_6(L)$ の法線微分が階数二を持てず、
$Z$ の被約性に反する。これが曲線による二十一点の分離を与える。
\end{proof}


\section{滑らかなGram曲線と点配置に付随する二重被覆}
\label{sec:gram-prym-v25}

命題\ref{prop:prym-baseline-v26}の行列式曲線を、線束を伴う滑らかな曲線へ改善する。
全自己反対incidenceと指定反対合の枠を固定する。
対称行列表示と二捩れ線束の対応は既知の理論である
\cite[\S2, Proposition 4.6]{BeauvilleDeterminantal}。
ここで証明する本族に固有の点は、既存の固定成分内で滑らかさが実現することと、
その線束を二十一点配置の非固定部分から復元できることである。

\begin{theorem}[滑らかな四次・六次曲線の実在]
\label{thm:gram-smooth-v25}
$Z_3$ と $Z_1$ は、それぞれ、命題\ref{prop:prym-baseline-v26}の被約二十一点配置の条件に加えて
\[
 C=\Gamma_P=V(\det(PAP^t))
\]
が滑らかな平面四次曲線、滑らかな平面六次曲線となる非空開部分を持つ。
従ってこれらの曲線は幾何学的に既約であり、種数はそれぞれ三、十である。
$Z_2$ の非退化二次曲線と合わせ、三族とも滑らかなGram曲線を持つ開部分が存在する。
\end{theorem}

\begin{proof}
命題\ref{prop:prym-baseline-v26}と同じ素数 $101$、方向選択seed、基点、pivot変数、自由変数を用いる。
固定成分 $Z_3,Z_1$ のHensel解を精度 $101^{256}$ まで再構成する。
自由変数の全リストを既存の被約二十一点証明書と照合しているので、
これらは命題\ref{prop:prym-baseline-v26}で同一成分への所属と被約二十一点性を証明した、
同じ代数的 $101$ 進点の有限精度表示である。

$F=\det(PAP^t)$、$d=\deg F=2p$ と置く。
$F_x,F_y,F_z$ の全てに次数 $2d-4$ の単項式を掛けたものを、
次数 $3d-5$ の単項式基底で並べた行列を $\mathcal M_F$ とする。
次の小行列式を証明書で指定する。
\[
\begin{array}{c|c|c|c|c|c}
 \text{成分}&d&3d-5&\mathcal M_F\text{ の大きさ}&
 \text{指定小行列式の付値}&\text{保証精度}\\\hline
 Z_3&4&7&45\times36&11&256\\
 Z_1&6&13&135\times105&57&256
\end{array}
\]
前者は三十六次、後者は百五次の小行列式である。
$F$ の係数、偏微分、Macaulay行列、小行列式はいずれも
Hensel座標の整係数多項式であり、分母を導入しない。
従って入力の $101^{256}$ 精度が小行列式にもそのまま保証される。
付値が保証精度未満なので、真のHensel解で小行列式は非零である。

よって $F_x,F_y,F_z$ が生成する斉次イデアルは
次数 $3d-5$ の全ての単項式を含む。
もし代数閉包上の射影点で三偏微分が同時に零になれば、
その点では同次数の全ての単項式が零となり矛盾する。
従って射影的な共通零点はなく、$C$ は幾何学的に滑らかである。
正標数での退化した基点に滑らかさを主張しているのではなく、
標数零のHensel点における非消滅を証明していることに注意する。

滑らかさは開条件なので、同じ既存成分上に求める非空開部分を得る。
滑らかな正次数平面曲線は、異なる正次数成分があれば交点を持って特異となるため既約である。
種数は平面曲線の随伴公式から $(d-1)(d-2)/2$ である。
\end{proof}

\paragraph{再現入口.}
\nolinkurl{gram_curves_v25.py} と
\nolinkurl{gram_curves_v25_certificate.json} が指定小行列式の行番号、
正規化した行列式の剰余、Gram多項式の係数hash、既存証明書との同一性を記録する。
通常の実行は読取り専用である。
この検証は滑らかさのための追加の有限証明書であり、
同一成分への所属と被約二十一点性の証明は命題\ref{prop:prym-baseline-v26}のものを使用する。

\begin{proposition}[非自明二捩れ線束と非固定点の除数]
\label{prop:gram-prym-divisor-v25}
命題\ref{prop:prym-baseline-v26}の一般署名命題の仮定に加えて、$C=\Gamma_P$ が滑らかであるとする。
$p=1,2,3$、$q=6-p$ とし、$i:C\hookrightarrow\mathbb P^2$ と書く。
対称行列 $Q_P=PAP^t$ の余核を
\begin{equation}\label{eq:gram-line-resolution-v25}
 0\longrightarrow\mathcal O_{\mathbb P^2}(-2)^p
 \xrightarrow{Q_P}\mathcal O_{\mathbb P^2}^{p}
 \longrightarrow i_*\mathcal L\longrightarrow0
\end{equation}
で定め、$\eta=\mathcal L(1-p)$ と置く。
このとき $\mathcal L$ は $C$ 上の線束であり、
\[
 \eta^{\otimes2}\simeq\mathcal O_C.
\]
$p=2,3$ では $H^0(C,\eta)=0$ なので $\eta$ は非自明である。
$p=1$ では $\eta=\mathcal O_C$ となる。

非固定点を被約有効除数
\[
 D_{\rm mov}=\sum_{x\in Z\setminus\operatorname{Fix}(\iota)}x
\]
と見なすと、
\begin{equation}\label{eq:moving-prym-divisor-v25}
 \mathcal O_C(D_{\rm mov})\simeq\eta(q),
 \qquad
 \eta\simeq\mathcal O_C(D_{\rm mov}-qH)
\end{equation}
となる。従ってこの二捩れ線束は、指定反対合に付随するGram曲線と
非固定点の除数からも復元できる。
\end{proposition}

\begin{proof}
$C$ 上で $Q_P$ の階数は常に $p-1$ である。
実際、階数が $p-2$ 以下なら随伴行列が零で、
$\det Q_P$ の全微分が零となり $C$ の滑らかさに反する。
局所的に単元ブロックを消すと $Q_P$ は
$\operatorname{diag}(I_{p-1},f)$ と同値になり、$f$ は $C$ の局所方程式である。
従って余核は $C$ 上階数一の局所自由加群である。

式\eqref{eq:gram-line-resolution-v25}を
$\mathcal H om(-,\mathcal O_{\mathbb P^2}(-2))$ で双対化する。
$Q_P$ の対称性により、得られる余核も $i_*\mathcal L$ である。
一方、次数 $2p$ のCartier因子の双対性から
\[
 \mathcal E xt^1_{\mathbb P^2}
 (i_*\mathcal L,\mathcal O_{\mathbb P^2}(-2))
 \simeq i_*\bigl(\mathcal L^{-1}(2p-2)\bigr).
\]
よって $\mathcal L^{\otimes2}\simeq\mathcal O_C(2p-2)$、
従って $\eta^{\otimes2}\simeq\mathcal O_C$ である。
分解を $1-p$ だけねじると
\[
 0\longrightarrow\mathcal O(-p-1)^p
 \longrightarrow\mathcal O(1-p)^p
 \longrightarrow i_*\eta\longrightarrow0.
\]
$p=2,3$ では $H^0(\mathcal O(1-p))=0$ と
$H^1(\mathbb P^2,\mathcal O(-p-1))=0$ により $H^0(C,\eta)=0$ を得る。
$p=1$ は一行一列の表示そのものである。

次に $C$ 上の自明束 $\mathcal V=\mathcal O_C^7$ に $A$ の非退化対称形式を入れ、
$P$ の行空間を部分束
\[
 \mathcal S\simeq\mathcal O_C(-1)^p\subset\mathcal V
\]
とする。その根基 $\mathcal T=\mathcal S\cap\mathcal S^\perp$ は線束である。
式\eqref{eq:gram-line-resolution-v25}を $C$ に制限した際のTor項を計算し、
一つねじると
\[
 \mathcal T\simeq\mathcal L(1-2p)=\eta(-p)
\]
を得る。

命題\ref{prop:prym-baseline-v26}により $B$ の階数低下点は固定点に一致し、固定点は $C$ の外にある。
従って $C$ 上では $B^t$ の像も部分束
\[
 \mathcal B\simeq\mathcal O_C(-1)^q\subset\mathcal S^\perp
\]
となる。直交性は全incidenceの混合ブロックから従う。
商 $\mathcal F=\mathcal S^\perp/\mathcal B$ は線束であり、
\[
 \det(\mathcal S^\perp)\simeq\mathcal O_C(-p),\qquad
 \mathcal F\simeq\mathcal O_C(q-p)
\]
である。従って自然な射
\[
 \mathcal T\longrightarrow\mathcal F
\]
は $\eta^{-1}(q)\simeq\eta(q)$ の切断を定める。
その零点は $\mathcal T(x)\subset\mathcal B(x)$ となる点である。
これは $\mathcal S(x)\cap\mathcal B(x)\ne0$、すなわち
$\operatorname{rank}L(x)<6$ と同値であり、$C$ 上では非固定点ちょうど全部である。
点スキームが有限なので切断は恒等的には零でない。
この切断の零除数の次数は
\[
 \deg\eta(q)=2pq
\]
で、既に確定した相異なる非固定点数も $2pq$ である。
従って各零点の重複度は一であり、式\eqref{eq:moving-prym-divisor-v25}を得る。
\end{proof}

\begin{corollary}[Prym曲線とPrym-canonical除数]
\label{cor:gram-prym-canonical-v25}
上の非空開部分で、$Z_3$ から種数三の曲線と非自明二捩れ線束、
$Z_1$ から種数十の曲線と非自明二捩れ線束を得る。
対応する連結エタール二重被覆の種数はそれぞれ五、十九である。
特に $Z_1$ では
\[
 K_C\simeq\mathcal O_C(3),\qquad
 D_{\rm mov}\in|K_C\otimes\eta|,\qquad
 h^0(C,K_C\otimes\eta)=9.
\]
従って十八個の非固定点は、Prym-canonical線形系の被約除数をなす。
\end{corollary}

\begin{proof}
$\eta^{\otimes2}\simeq\mathcal O_C$ の同型を選ぶと
$\operatorname{Spec}_C(\mathcal O_C\oplus\eta)$ はエタール二重被覆となる。
$\eta$ が非自明なので連結であり、Riemann--Hurwitzから種数は $2g(C)-1$ である。
六次曲線では随伴公式により $K_C=\mathcal O_C(3)$ であり、$q=3$ なので
式\eqref{eq:moving-prym-divisor-v25}を適用できる。
Riemann--Rochと $h^0(C,\eta)=0$ により
$h^0(C,K_C\otimes\eta)=g(C)-1=9$ となる。
\end{proof}

\begin{proposition}[補助Gram曲線との二重交差]
\label{prop:complementary-gram-contact-v25}
上と同じ仮定の下で
\[
 \Gamma_B=V(\det(BAB^t))
\]
は次数 $2q$ の非零行列式曲線であり、$C$ を成分に持たない。
$C$ 上のCartier除数として
\begin{equation}\label{eq:complementary-gram-contact-v25}
 \Gamma_B|_C=2D_{\rm mov}
\end{equation}
が成立する。特に $Z_1$ では二つの六次行列式曲線の交差は、
十八個の非固定点に重複度二ずつ集中する。
また、全ての固定点は $\Gamma_B$ の特異点である。
補助曲線の既約性と特異点は第\ref{sec:complementary-nodes-v26}節で決定する。
\end{proposition}

\begin{proof}
$\mathcal S^\perp$ の根基も $\mathcal T$ なので、
\[
 \mathcal H=\mathcal S^\perp/\mathcal T
\]
は階数 $q$ の非退化対称束であり、$\det\mathcal H\simeq\eta^{-1}$ である。
$\mathcal B\to\mathcal H$ は同階数間の射で、その行列式は
$\eta^{-1}(q)$ の切断である。
これが退化するのは $\mathcal T\cap\mathcal B\ne0$ の点なので、
前の命題と同じ次数計算により、その零除数は $D_{\rm mov}$ である。
$\mathcal B$ 上のGram形式は、$\mathcal H$ の非退化形式をこの射で引き戻したものである。
従ってその行列式の零除数は二倍され、
式\eqref{eq:complementary-gram-contact-v25}を得る。
非固定点は有限なので行列式の $C$ への制限は恒等的には零でない。
これにより多項式自体が非零であり、$C$ を成分に持たないことも分かる。

固定点では $B$ の階数は $q-1$ である。
$G=BAB^t$ の階数が $q-2$ 以下なら、$\operatorname{adj}(G)=0$ なので
$d(\det G)=0$ である。
階数が $q-1$ なら、その核の生成元 $v$ は $B^tv=0$ を満たし、
\[
 v^t(dG)v=v^t(dB)AB^tv+v^tBA(dB)^tv=0.
\]
随伴行列は $vv^t$ の非零倍なので、この場合にも $d(\det G)=0$ となる。
従って固定点は $\Gamma_B$ の特異点である。
\end{proof}

\begin{remark}[二つの二重構造の区別]
$\eta$ が定めるエタール二重被覆は、曲線 $C$ の上の被覆である。
一方、指定反対合は $D_{\rm mov}$ の点を二点ずつ交換する。
この二つを同一視する主張はしていない。
$p=3$ の対称 $3\times3$ 二次式行列 $Q_P$ は、
$\mathbb P^2\times\mathbb P^2$ の双次数 $(2,2)$ 超曲面
$y^tQ_P(x)y=0$ も定め、その第一射影のconic bundleの判別式は $C$ である。
そのPrymデータは上の $\eta$ で記述される既知の構成に属する
\cite[\S1--2]{BoehningBothmerPrym}。
本族が一般の双次数 $(2,2)$ 超曲面を支配するとは主張せず、
一般のconic bundleの非安定有理性をこの特別な部分族へ転用しない。
\end{remark}

\section{六次Gram曲線からの復元}
\label{sec:prym-torelli-sextic-v26}

署名 $(3,3)$ の成分 $Z_1$ では、二つの固有空間の順序を保つ。
すなわち $P$ に用いる三次元固有空間を指定し、
関係基底変換、二つのブロック内の基底変換、外側の座標変換による有理商を
$\mathfrak Z_1^+$ と書く。
その次元は命題\ref{prop:prym-baseline-v26}と同じく七である。
まずこの順序を保持して逆構成を記述する。
第\ref{subsec:unordered-blocks-v27}項では、順序を忘れた商も同じ関数体を持つことを示す。

\begin{theorem}[平面Prymデータによる復元]
\label{thm:plane-prym-torelli-v26}
$\mathfrak Z_1^+$ の非空開部分では、平面Prymデータ
\[
 (C,\eta,H),\qquad
 C=V(\det(PAP^t)),\quad H=\mathcal O_C(1),\quad
 \eta=\operatorname{coker}(PAP^t)(-2)
\]
が、指定した順序を保つ同型を除いて対象を一意に復元する。
従ってこの写像は像への双有理写像であり、その像の次元は七である。
平面直線束の標識は、第\ref{subsec:abstract-prym-torelli-v27}項で除く。
指定反対合そのものの標識は保持する。
\end{theorem}

\begin{proof}
まず、既存の $Z_1$ の非空開部分で次の二つの階数条件が成立することを確認する。
$P(x)=\sum_{i=0}^2x_iP_i$ の九つの係数行ベクトルを並べた
$9\times7$ 行列を $\widehat P$ とする。
また
\[
 K_{P,A}=\{b(x)\in k^7\otimes H^0(\mathcal O(1)):
                              P(x)A b(x)^t=0\}
\]
と置く。
第\ref{sec:gram-prym-v25}節と同じ素数 $101$、seed $17029$ のHensel点で、
以下の小行列式はいずれも $101$ 進単元である。
\[
\begin{array}{c|c|c}
 \text{行列}&\text{階数}&\text{指定小行列式の }101^3\text{ における剰余}\\\hline
 \widehat P\ (9\times7)&7&755782\\
 \widehat P A\widehat P^t\text{ の指定主 }7\times7\text{ 部分行列}&7&389051\\
 PAb^t=0\text{ の係数行列 }(18\times21)&18&990506\\
 B\text{ の三行の }K_{P,A}\text{ 内での座標行列}&3&60902
\end{array}
\]
従って $\operatorname{rank}\widehat P=7$、$\dim K_{P,A}=3$、
そして $B$ の三行は $K_{P,A}$ 全体の基底をなす。
これは $101$ 進点だけの性質ではなく、同じ既約成分上の非空開条件である。
この開部分を、二十一点の被約性と $C$ の滑らかさの既存の非空開部分と交わらせる。

次に $Q=PAP^t$ から $(P,A)$ を復元する。
$Q_{uv}$ の $x_i x_j$ の係数を $q_{uv,ij}$ とし、
九行九列の行列 $G(Q)$ を
\begin{equation}\label{eq:sextic-coefficient-gram-v26}
 G(Q)_{(u,i),(v,j)}=
 \begin{cases}
  q_{uv,ii},&i=j,\\
  q_{uv,ij}/2,&i\ne j
 \end{cases}
 \qquad (0\le u,v,i,j\le2)
\end{equation}
で定める。添字を入れ替える場合には $Q_{vu}=Q_{uv}$ とする。
同符号ブロックの全incidenceは
\[
 (P_i)_u A(P_j)_v^t=(P_j)_u A(P_i)_v^t
\]
を含むため、式\eqref{eq:sextic-coefficient-gram-v26}は
\[
 G(Q)=\widehat P A\widehat P^t
\]
を正確に復元する。
特にその階数は七である。

可逆な主部分行列 $G_I=G(Q)_{I,I}$ を選び、
\begin{equation}\label{eq:sextic-rational-factorization-v26}
 A_0=G_I,\qquad
 \widehat P_0=G(Q)_{*,I}G_I^{-1}
\end{equation}
と置く。
階数七のSchur補行列が零になることから
$G(Q)=\widehat P_0A_0\widehat P_0^t$ である。
これは $Q$ の係数の有理式による復元である。
他の非退化な階数七の分解
$G(Q)=\widehat P_1A_1\widehat P_1^t$ があれば、
$\operatorname{im}\widehat P_1=\operatorname{im}G(Q)
=\operatorname{im}\widehat P_0$ なので
\[
 \widehat P_1=\widehat P_0g,\qquad
 A_1=g^{-1}A_0g^{-t}
 \quad(g\in\operatorname{GL}_7)
\]
となる。従って $(P,A)$ は関係基底変換だけの差を除いて一意である。

復元された $(P,A)$ から十八行二十一列の線形系
$PAb^t=0$ を作る。その核は三次元であり、$B$ の三行はこの核全体を張る。
従って $B$ も $\operatorname{GL}_3$ による行変換だけの差を除いて一意である。
ここで $Q$ は元の成分の像に属するものを扱っているので、
復元した核の基底が満たす $B$--$B$ ブロックのincidenceは元の対象から従う。
任意の階数七の $G(Q)$ についてそれらの方程式が自動的に成立するとは主張しない。

最後に平面Prymデータから $Q$ を復元する。
このデータは $\eta(2H)$ とその最小次数付き自由分解
\[
 0\longrightarrow\mathcal O_{\mathbb P^2}(-2)^3
 \xrightarrow{Q}\mathcal O_{\mathbb P^2}^{3}
 \longrightarrow i_*\eta(2H)\longrightarrow0
\]
を同型を除いて決める。
対称な最小分解の一意性により、二つの対称表示は合同変換で結ばれる
\cite[\S2.3]{BeauvilleDeterminantal}。
従って $(C,\eta,H)$ は $Q$ を正ブロックの $\operatorname{GL}_3$ と
平面の $\operatorname{PGL}_3$ の作用だけの差を除いて決める。
前二段階の復元と合わせて、$\mathfrak Z_1^+$ の一般点は一意に決まる。
標数零の既約な粗有理商の間で、一般の幾何学的ファイバーが一点なので、
像への写像は双有理である。像の次元は従って七である。

\end{proof}

\paragraph{Gram像の次元の独立検算.}
同じHensel点の正規化チャートは百五十三変数、七十二方程式を持ち、接次元は八十一である。
三十六個のGram係数の微分をこの接空間へ制限すると、指定二十四次小行列式が単元となる。
その $101^3$ における剰余は $822512$ であり、従って像の次元は少なくとも二十四である。
一方、正規化前の全pairingチャートは八十二次元であり、
関係基底の $\operatorname{GL}_7$ と負ブロックの $\operatorname{GL}_3$ は
$Q$ を変えない五十八次元の自由な作用を持つ。
自由性は $\operatorname{rank}\widehat P=7$ と $B$ の行独立性から従う。
従ってGram像の次元は高々 $82-49-9=24$ であり、正確に二十四である。
正ブロックと座標の二つの $\operatorname{GL}_3$ の作用は、
互いに逆のスカラーによる一次元核を持ち、有効軌道の次元は十七である。
従って商の次元は $24-17=7$ となり、上の復元による次元計算と一致する。

\paragraph{再現入口.}
\nolinkurl{gram_dominance_v26.py} を引数 \texttt{--signature 3} で実行する。
対応する \nolinkurl{gram_dominance_v26_p3_certificate.json} は、
階数証明書、線形核の基底、Gram分解の有理式の照合、
二十四次微分小行列式を含む。
全ての入力精度は $101^{32}$ であり、ここで使う小行列式は全て単元である。
従って階数と復元の成立は既存の同じHensel点で検証されている。

\subsection{平面表示の内在性と抽象Prymデータからの復元}
\label{subsec:abstract-prym-torelli-v27}

平面曲線の直線束の一意性は古典的なNoether理論の一部である
\cite{HartshorneNoether,CoppensKatoPlaneSystems}。
ここで必要な場合は、随伴公式と点の分離だけで証明できる。
以下の補題を用いると、六次の場合の復元から平面直線束の標識を除ける。

\begin{lemma}[平面直線束の一意性と剛性]
\label{lem:intrinsic-plane-series-v27}
標数零の代数閉体上で、$C\subset\mathbb P^2$ を滑らかな次数 $d\ge4$ の曲線とし、
$H=\mathcal O_C(1)$ と置く。
次数が $d$ 以下で $h^0(C,M)\ge3$ を満たす線束 $M$ は、$H$ だけである。
さらに、通常の決定因子的スキーム構造を入れたBrill--Noether軌跡は
\[
 W_d^2(C)=\{[H]\}
\]
という被約な一点である。
従って、次数 $d$ の二つの滑らかな平面曲線の任意の同型は、平面の射影変換へ一意に延長する。
\end{lemma}

\begin{proof}
$M$ の固定因子を除いた線束を $M'$、その次数を $n\le d$ とする。
$h^0(M')=h^0(M)\ge3$ であり、標数零のBertini定理により、
$|M'|$ の一般元 $D=p_1+\cdots+p_n$ は相異なる点からなる。
随伴公式は $K_C=(d-3)H$ を与え、
\[
 H^0(\mathbb P^2,\mathcal O(d-3))\xrightarrow{\sim}H^0(C,K_C)
\]
である。Riemann--Rochにより、$D$ におけるこの空間の評価写像の階数は
\begin{equation}\label{eq:plane-series-evaluation-v27}
 \operatorname{rank}(\mathrm{ev}_D)
 =h^0(K_C)-h^0(K_C-D)
 =n+1-h^0(M')\le n-2
\end{equation}
となる。

$n\le d-2$ なら、各 $p_i$ を除く $n-1$ 点を通り $p_i$ を避ける直線を一つずつ選ぶ。
その積を、必要なら $p_i$ を避ける一次式で次数 $d-3$ まで補えば、
$p_i$ 以外で零、$p_i$ で非零となる形式を得る。
従って評価の階数は $n$ であり、\eqref{eq:plane-series-evaluation-v27}に反する。

次に、$d-1$ 個の相異なる非共線点は次数 $d-3$ の形式に独立な条件を課す。
実際、そのうちの一点 $p$ を固定すると、残りの点には、結ぶ直線が $p$ を通らない二点がある。
その直線と、残る $d-4$ 点をそれぞれ通り $p$ を避ける直線の積が $p$ を分離する。
この議論は $d=4$ の場合にも、三非共線点を一次式で分離する議論として成立する。

$n=d-1$ で $D$ が非共線なら、今示したことにより階数は $d-1$ となる。
$D$ が共線でも、直線への制限は次数 $d-3$ の多項式全体を与えるので、階数は $d-2$ である。
どちらも\eqref{eq:plane-series-evaluation-v27}の上界 $d-3$ に反する。
従って $n=d$ である。
このとき $D$ が非共線なら、$d-1$ 個の非共線部分集合を選べ、評価階数は少なくとも $d-1$ となる。
これも上界 $d-2$ に反するため、$D$ は一本の直線上にある。
B\'ezoutの定理から $D$ はその直線と $C$ の交わりそのものであり、$M'=H$ である。
$\deg M\le d=n$ なので固定因子は零、従って $M=H$ となる。

被約性を確認する。$h^0(H)=3$ なので、$[H]$ における $W_d^2(C)$ の接空間は、
三つの切断が全て持ち上がる一次変形 $\xi\in H^1(C,\mathcal O_C)$ からなる。
その条件は、cup積
\[
 H^0(C,H)\xrightarrow{\ \xi\ }H^1(C,H)
\]
が零になることである。Serre双対性により、その接空間は乗法写像
\[
 H^0(C,H)\otimes H^0(C,K_C-H)\longrightarrow H^0(C,K_C)
\]
の像の零化空間である。この乗法は平面上の
\[
 S_1\otimes S_{d-4}\longrightarrow S_{d-3},
 \qquad S=k[x_0,x_1,x_2]
\]
と一致し、全射である。従って接空間は零である。
既に閉点が $[H]$ 一点と分かっているので、Nakayamaの補題によりその局所環は体である。

最後に、二つの平面曲線の同型はそれぞれの直線束を対応させる。
三次元の切断空間の同型を射影化すると平面の射影変換が得られる。
曲線は平面を張るので、その延長は一意である。
\end{proof}

\begin{theorem}[抽象Prymデータによる六次の復元]
\label{thm:abstract-sextic-prym-torelli-v27}
署名 $(3,3)$ の七次元の標識付き有理商について、写像
\[
 \Phi_1:\mathfrak Z_1^+\dashrightarrow\mathcal R_{10},\qquad
 [L,A]\longmapsto(C,\eta)
\]
は、その七次元の像への双有理写像である。
従って一般の抽象Prym対 $(C,\eta)$ から、平面表示、関係代数、全二十一点配置、
および特定の十八点除数 $D_{\mathrm{mov}}\in|K_C\otimes\eta|$ を復元できる。
\end{theorem}

\begin{proof}
この像に属する $C$ は滑らかな平面六次曲線である。
補題\ref{lem:intrinsic-plane-series-v27}により、$H$ は
$W_6^2(C)$ の唯一の被約な点として $(C,\eta)$ から内在的に決まる。
二つの抽象Prym対の同型は $H$ も保つので、平面の射影変換へ延長する。
従って定理\ref{thm:plane-prym-torelli-v26}の平面Prym復元を適用でき、
一般の幾何学的ファイバーは一点となる。
標数零なので、既約な有理商からその像への写像は双有理である。
全点配置と $D_{\mathrm{mov}}$ の復元は、同定理のGram分解と混合syzygy核の逆構成に従う。
\end{proof}

\paragraph{線束の三乗根との違い.}
随伴公式 $H^{\otimes3}=K_C$ だけを使えば、可能性が有限個であることしか分からない。
ここでは切断数 $h^0(H)=3$ を併せて使い、実際の平面表示が一意であることを示している。
この議論は曲線と二捩れ線束を保持しており、Prym多様体だけからの復元を主張しない。

\subsection{二つの固有ブロックの順序を忘れる}
\label{subsec:unordered-blocks-v27}

反対合の標識を保持したまま、三次元の二固有ブロックの順序だけを忘れることもできる。
以下では、この忘却を反対合そのものの忘却と区別して定義する。
pairingのスカラー正規化を行う前の既存の枠付き成分の非空開部分 $\mathcal U$ を、
$P$ が平面全体で階数三、$B$ が指定した三固定点で階数低下し、
上の全復元条件を満たすように取る。
$\mathcal G_{\mathrm{ord}}$ は関係枠、二ブロック内の行基底、平面座標、およびpairingのスカラーを変える群とする。
ブロック交換を
\[
 s(P,B,A)=(B,P,A)
\]
で定める。$sJs^{-1}=-J$ であり、全incidenceは $J$ の符号を変えても同じなので、
$s$ は全incidence空間の自己同型である。

\begin{proposition}[固有ブロックの内在的な選択]
\label{prop:unordered-blocks-v27}
上の開部分で、二つの固有ブロックのうち平面全体で階数三を持つものは一つだけである。
従って、順序を忘れた商を
\[
 \mathfrak Z_1^{\mathrm{unord}}
 :=(\mathcal U\sqcup s\mathcal U)/
   (\mathcal G_{\mathrm{ord}}\rtimes\langle s\rangle)
\]
の有理商として定義すると、自然写像
$\mathfrak Z_1^+\dashrightarrow\mathfrak Z_1^{\mathrm{unord}}$ は双有理である。
写像 $\Phi_1$ は、この商上でも像への双有理写像を与える。
\end{proposition}

\begin{proof}
$P$ は全点で階数三を持つ一方、$B$ は三固定点で階数低下するので、
$\mathcal U\cap s\mathcal U=\varnothing$ である。
二ブロックのうち全点で最大階数を持つものを選べば、順序は一意に戻る。
群 $\mathcal G_{\mathrm{ord}}$ はこの条件を保ち、$s$ は二つの開部分を交換する。
従って二つの関数体の積の不変体は
\[
 \bigl(k(\mathcal U)\times k(s\mathcal U)\bigr)^{{\mathcal G}_{\mathrm{ord}}\rtimes\langle s\rangle}
 \simeq k(\mathcal U)^{{\mathcal G}_{\mathrm{ord}}}
\]
である。これが商の双有理同定を与える。
全点で最大階数を持つブロックを用いて $(C,\eta)$ を作れば、
その構成は順序を忘れた商上でも定義され、前定理を適用できる。
\end{proof}

\paragraph{標識の適用範囲.}
ここで除いた標識は、固定された射影的反対合の二固有ブロックの順序だけである。
一つの代数上に異なる反対合が存在するか、その全てが共役か、一般の全外部自己同型群が自明かは、
この議論では決定しない。

\section{六次Prym像の方程式と標準形}
\label{sec:intrinsic-sextic-image-v28}

前節の復元は、像に属するPrym対から元の対象を取り戻すものであった。
ここでは、その像を判定する追加条件を九本の方程式として与える。
併せて、十八点除数を直接定める標準的なPrym-canonical切断を構成する。
以下の標準形は、既存の成分 $Z_1$ と交わる非空開部分で用いる。
九本の方程式の全零点集合について、既約性や有理性を主張するものではない。

三次元空間を $U,V$ とし、対称な三行三列二次式行列 $Q$ の係数Gram形式
$G(Q)$ を式\eqref{eq:sextic-coefficient-gram-v26}で定める。
これは $U\otimes V$ 上の対称形式である。
$\operatorname{rank}G(Q)=7$ とし、その二次元の根基を $K$ と書く。
$K$ の元を三行三列行列と見たとき、行列式が $\mathbb P(K)$ 上で
相異なる三点に零を持つ場合を、\emph{正則な核pencil}と呼ぶ。
これは二つの空間の基底と $K$ の基底に依存しない開条件である。

\begin{theorem}[九本の余因子方程式]
\label{thm:sextic-normal-form-v28}
係数Gram形式の階数が七で、核pencilが正則であるとする。
行基底と平面座標を変換すると、$Q=PAP^t$ を
\begin{equation}\label{eq:canonical-P-R-v28}
\begin{split}
 P&=\begin{pmatrix}
 x_0&x_1&x_2&0&0&0&0\\
 x_1&0&0&x_0&x_2&0&0\\
 x_2&0&0&0&0&x_0&x_1
 \end{pmatrix},\\[2pt]
 R&=\begin{pmatrix}
 0&x_2&-x_1&0&0&0&0\\
 0&0&0&x_2&-x_0&0&0\\
 0&0&0&0&0&x_1&-x_0
 \end{pmatrix}
\end{split}
\end{equation}
および次の十九変数の対称行列で表せる。
\begin{equation}\label{eq:A-nineteen-v28}
 A=\begin{pmatrix}
 s&f&g&l&m&p&r\\
 f&a&b&s&g&r&e\\
 g&b&c&m&d&s&f\\
 l&s&m&h&i&k&p\\
 m&g&d&i&j&l&s\\
 p&r&s&k&l&n&o\\
 r&e&f&p&s&o&q
 \end{pmatrix},\qquad \det A\ne0.
\end{equation}
添字を $0,\ldots,6$ とし、
\[
 \mathcal J=\{(1,4),(1,5),(1,6),(2,3),(2,4),(2,6),
                    (3,5),(3,6),(4,5)\}
\]
と置く。このとき、混合直交syzygyを全部用いた負ブロックは
\[
 B=RA^{-1}
\]
であり、$(P,B,A)$ が全自己反対incidenceを満たす必要十分条件は
\begin{equation}\label{eq:nine-cofactor-equations-v28}
 (A^{-1})_{uv}=0\quad ((u,v)\in\mathcal J),
 \qquad\text{すなわち}\qquad
 \operatorname{adj}(A)_{uv}=0\quad ((u,v)\in\mathcal J)
\end{equation}
である。

$s=1$、他の十八変数が零の点を $A_0$ とする。
この点は既存の整数基点 $B_-$ の標準形であり、
九本の余因子方程式の零点集合は $A_0$ で滑らかな十次元多様体である。
$A_0$ を通る既約成分を $\mathcal Y_0$ とすると、共通の非空開部分上で
\begin{equation}\label{eq:normal-form-quotient-v28}
 \mathfrak Z_1^+\sim_{\rm bir}
 \mathcal Y_0/(\mathbb G_m^3\rtimes\mathfrak S_3).
\end{equation}
従って、核pencilを正則とする階数七のGram軌跡の中で、
本研究の六次Prym像の枠付き成分は一般点で余次元九の局所完全交叉である。
その次元は二十四、平面座標と対称表示の基底を除いた像の次元は七である。
\end{theorem}

\begin{proof}
まず核pencilを標準化する。
$K$ の可逆な元を一つ選んで単位行列とし、もう一つの元を通常の行列と見なす。
行列式の三根が相異なるので、この行列は三つの相異なる固有値を持ち、対角化できる。
従って $K$ は三次元の対角空間の中の二次元部分空間となる。
この対角空間の各座標軸は $K$ に含まれない。
実際、座標軸が含まれれば、その階数一の元において行列式の零点が重複する。
従って各対角座標を非零スカラー倍して
\[
 K=\langle e_{00}-e_{11},\ e_{11}-e_{22}\rangle
\]
とできる。商 $(U\otimes V)/K$ の基底を
\[
 e_0=\overline e_{00}=\overline e_{11}=\overline e_{22},\quad
 e_1=\overline e_{01},\ e_2=\overline e_{02},\
 e_3=\overline e_{10},\ e_4=\overline e_{12},\
 e_5=\overline e_{20},\ e_6=\overline e_{21}
\]
とすれば、係数写像は\eqref{eq:canonical-P-R-v28}の $P$ となる。
$P$ の三次小行列式には符号を除いて $x_0^3,x_1^3,x_2^3$ が含まれるため、
$P$ は平面全体で階数三を持つ。
$G(Q)$ はこの商に非退化対称形式 $A$ を誘導する。
係数Gram形式の二組の添字に関する対称性は、
$P_iAP_j^t=P_jAP_i^t$ という九本の独立な線形条件と同値である。
この線形条件を解いた一般解が\eqref{eq:A-nineteen-v28}である。

次に一次式行ベクトル $b$ に対し、$c=Ab^t$ と置く。
$P c=0$ の十八係数を比較すると、$c$ の第零成分は零であり、残る成分は
\[
 \tau_0(0,x_2,-x_1,0,0,0,0)^t+
 \tau_1(0,0,0,x_2,-x_0,0,0)^t+
 \tau_2(0,0,0,0,0,x_1,-x_0)^t
\]
に限られる。従って混合syzygy空間は正確に三次元であり、
$R$ の三行がその基底を与える。
負ブロックの行基底を変える自由を除けば $B=RA^{-1}$ である。
このとき $PAB^t=PR^t=0$ は恒等的に成り立つ。

$A^{-1}$ を一時的に任意の対称行列 $Z$ と置く。
$B_iAB_j^t=B_jAB_i^t$ は $R_iZR_j^t=R_jZR_i^t$ と同値である。
\eqref{eq:canonical-P-R-v28}を代入すると、この九条件は符号を除いて
$Z_{uv}=0$、$(u,v)\in\mathcal J$ そのものである。
これで全incidenceについての必要十分性を得る。

$A_0$ は $e_0$ を固定し、$(e_1,e_3)$、$(e_2,e_5)$、$(e_4,e_6)$ を交換する行列である。
従って $A_0^{-1}=A_0$、$\det A_0=-1$ である。
逆行列の微分 $d(A^{-1})=-A^{-1}(dA)A^{-1}$ を用いると、
$\mathcal J$ の上記の順序で並べた九本の余因子の
\[
 (p,m,i,r,o,l,b,g,f)
\]
に関するJacobi行列は $A_0$ で単位行列となる。
従って $A_0$ は滑らかな十次元点であり、残りの
\[
 (s-1,a,c,d,e,h,j,k,n,q)
\]
は完備局所環の正則パラメータとなる。
既存の基点の平面座標を $(x_2,x_1,x_0)$ の順に並べ、
係数行の $0,1,2,3,5,6,7$ 番を関係基底として選ぶと、
その $P,A$ は上記の $P,A_0$ に一致する。
この変換は行列式が $\pm1$ の整数基底変換である。
特に核pencilの正則性は同じ既存成分上の非空開条件であり、
滑らかなGram六次曲線と被約二十一点の既存の開条件と同時に満たせる。

最後に商を記述する。
正則な二次元pencil全体は $\operatorname{GL}(U)\times\operatorname{GL}(V)$ の一つの軌道であり、
その次元は $\dim\operatorname{Gr}(2,9)=14$ である。
固定した核に対するGram形式は上の十九次元線形空間の開部分をなすので、
正則な核pencilを持つ階数七Gram軌跡は滑らかで既約な三十三次元空間である。
標準pencilの安定化群は、三つの固有方向を同時に置換する $\mathfrak S_3$ と、
\[
 \operatorname{diag}(a_0,a_1,a_2),\quad
 \operatorname{diag}(b_0,b_1,b_2),\qquad
 a_0b_0=a_1b_1=a_2b_2=\rho
\]
からなる四次元トーラスの半直積である。
実際、安定化群はpencilの三つの特異元を置換し、置換を除いた後には
両方の三つの固有方向を保つため、この対角形になる。
$(a_i,b_i)=(\lambda,\lambda^{-1})$ は全係数上で自明に作用する。
この一次元の核と残る有限スカラー核を除いてもトーラスは三次元であり、
そのパラメータを $\rho^2,a_0/a_2,a_1/a_2$ と取れば、
有効な安定化群は $\mathbb G_m^3\rtimes\mathfrak S_3$ となる。
$\mathfrak S_3$ は $A_0$ を固定し、トーラスは $A_0$ をその非零スカラー倍に移す。
$A_0$ を通る成分は唯一なので、$\mathcal Y_0$ はこの安定化群で保たれる。
標準形は $P,A$ と混合syzygyの全基底を復元するため、
この安定化群による商が元の枠を除いた商と一致する。
基点で選んだ成分について\eqref{eq:normal-form-quotient-v28}を得る。
枠付きGram成分の次元は $14+10=24$ である。
一般点の有効軌道次元十七は既存の復元開部分で確定しているので、商の次元は七となる。
\end{proof}

\paragraph{基底に依存しない判定.}
九条件は、標準化を行わずにも記述できる。
$W=(U\otimes V)/K$ に誘導される形式を $A$、係数評価を $P$ とする。
三次元の一次syzygy空間
\[
 T=\{b\in W\otimes V^*:P(x)Ab(x)^t=0\}
\]
上の自然なテンソル
\[
 \Omega\in\bigwedge^2T^*\otimes\bigwedge^2V^*,
 \qquad
 \Omega_{ij}(b,d)=b_iAd_j^t-b_jAd_i^t
\]
が零になることが、九条件に相当する。
従って $G(Q)$ の階数、核pencilの判別式、$\Omega=0$ はいずれも
平面Prym対から決まる条件である。
対称最小表示の存在と一意性を仮定し、必要な二十一点の被約性・右復元の開条件を加えれば、
これらは全自己反対incidenceに属することの必要十分な判定を与える。
本研究の像を選ぶ際には、さらに $A_0$ を通る成分 $\mathcal Y_0$ を指定する。
この成分の指定を省いて全零点集合を一つの既約像と同一視しない。

\subsection{核pencilが定める切断と十八点の最小分解}
\label{subsec:canonical-prym-section-v28}

標準形のベクトル $e_0$ が張る直線は、座標選択に依存しない。
これを記述するため、正則pencilの三つの階数二の元を $k_0,k_1,k_2$ とする。
各 $i$ について、他の二つの元の像の交わりを $U_i\subset U$ とし、
転置について同じ操作で $V_i\subset V$ を定める。
これらはそれぞれ一次元であり、
\[
 E=\bigoplus_{i=0}^2 U_i\otimes V_i\subset U\otimes V
\]
はpencilを含む三次元空間である。
この空間は三つの根の順序に依存しない。従って
\begin{equation}\label{eq:canonical-line-ell-v28}
 \ell=E/K\subset W
\end{equation}
は標準的な直線であり、標準形では $\ell=\langle e_0\rangle$ である。
また三つの $V_i$ が、平面 $\mathbb P(V)$ 内の標準的な三角形の頂点を定める。
標準形ではこれが三座標点 $T_\triangle$ である。

\begin{theorem}[標準的なPrym-canonical切断と十八点の方程式]
\label{thm:canonical-moving-section-v28}
定理\ref{thm:sextic-normal-form-v28}の設定で、
$C=V(\det Q)$ が滑らかで三角形 $T_\triangle$ を避けるとする。
\eqref{eq:canonical-line-ell-v28}の非零生成元 $v$ を選び、
\[
 t(x)=P(x)Av
\]
と置く。この一次式列ベクトルの余核への像は、非零切断
\[
 \bar t\in H^0(C,K_C\otimes\eta)
\]
をスカラーを除いて内在的に定める。
標準形で $v=e_0$ と取ると、この列ベクトルは
\[
 t(x)=\begin{pmatrix}
 sx_0+fx_1+gx_2\\ lx_0+sx_1+mx_2\\ px_0+rx_1+sx_2
 \end{pmatrix}
\]
である。
その零除数 $D$ は長さ十八で、飽和斉次イデアルと最小分解は
\begin{align}
 I_D&=I_3[\,Q\mid t\,]
      =(\det Q,\, (\operatorname{adj}(Q)t)_0,
                    (\operatorname{adj}(Q)t)_1,
                    (\operatorname{adj}(Q)t)_2),
 \label{eq:moving-ideal-v28}\\
 0&\longrightarrow S(-7)^3
 \longrightarrow S(-5)^3\oplus S(-6)
 \longrightarrow I_D\longrightarrow0,
 \qquad S=\mathbb C[x_0,x_1,x_2].
 \label{eq:moving-HB-v28}
\end{align}
さらに $L=\binom{P}{RA^{-1}}$ の最大小行列式が定める点スキームは
\[
 Z=D\sqcup T_\triangle
\]
であり、三角形の三点は被約である。
既存の被約二十一点開部分では $D=D_{\rm mov}$ である。
従って、七次元Prym像上の指定された十八点は、
負ブロックを計算することなく\eqref{eq:moving-ideal-v28}で復元できる。
\end{theorem}

\begin{proof}
この構成は商 $W$ の形式と標準的な直線 $\ell$ からなるため、
基底やpencilの三根の順序を変えても、切断の直線は変わらない。
$Q$ の対称表示の余核を $i_*\mathcal L$ とすると、
$\mathcal L=\eta(2H)$ であり、表示を一つねじって
\[
 0\longrightarrow\mathcal O(-1)^3
 \xrightarrow{Q}\mathcal O(1)^3
 \longrightarrow i_*(K_C\otimes\eta)\longrightarrow0
\]
を得る。従って $t$ は主張した切断を定める。
$P$ の全係数行は $W$ を張り、$A$ は非退化で $v\ne0$ なので、$t$ は零ではない。
$H^0(\mathcal O(-1))=H^1(\mathcal O(-1))=0$ により、$\bar t$ も零ではない。
$C$ は滑らかで既約なので、その零点は次数
$\deg(K_C\otimes\eta)=18$ の有効Cartier除数である。

$C$ 上では $Q$ の階数は二である。
局所的に二次の単元ブロックを消すと、$\bar t$ が零になる条件は
$[Q\mid t]$ の三次小行列式の消滅に一致する。
従って\eqref{eq:moving-ideal-v28}が $D$ を定める。
この最大小行列式イデアルは高さ二である。
Hilbert--Burch定理\cite[Theorem~2.2]{EisenbudSidman}を適用すると、
三本の五次生成元と一本の六次生成元、その間の三本の七次syzygyから
\eqref{eq:moving-HB-v28}を得る。
全ての行列要素が正次数なので分解は最小である。
高さ二の完全イデアルであるため、商は一次元Cohen--Macaulay環となり、イデアルは飽和している。
Hilbert級数からの独立な次数計算も
\[
 \deg D=\frac{3\cdot7^2-3\cdot5^2-6^2}{2}=18
\]
を与える。

残る点スキームとの同定を確認する。
標準形の $R$ は三角形の外で階数三であり、その核は
$P^t$ の三列と $e_0$ により生成される。
これら四列は三角形の外では独立である。
$L$ を右から $A$ 倍した行列は $\binom{PA}{R}$ なので、
$R$ の単元ブロックを消せば、六次小行列式イデアルは
$[PAP^t\mid PAe_0]=[Q\mid t]$ の三次小行列式イデアルになる。
従って三角形の外の点スキームは $D$ に等しい。

三角形の各頂点では $Q$ が可逆である。
したがって $P$ の行空間とその $A$ 直交補空間が局所的に直和をなし、
混合直交性を用いて $I_6(L)$ を $I_3(B)=I_3(R)$ に簡約できる。
$R$ の非零最大小行列式は符号を除き
\[
 x_0^2x_1,\ x_0^2x_2,\ x_0x_1^2,\ x_0x_1x_2,
 x_0x_2^2,\ x_1^2x_2,\ x_1x_2^2
\]
である。これは三角形のイデアル $(x_0x_1,x_0x_2,x_1x_2)$ に
斉次極大イデアルを掛けたものなので、射影平面上では被約な三頂点を定める。
$C$ が三角形を避けることから $D$ と三角形は交わらず、主張した非交和を得る。
\end{proof}

\begin{corollary}[五次式による十八点と三固定点の分離]
\label{cor:quintic-separation-v28}
前定理の設定で
\[
 h^0(\mathbb P^2,\mathcal I_D(5))=3,\qquad
 h^0(\mathbb P^2,\mathcal I_Z(5))=0.
\]
特に $D$ を通る五次式の空間は $\operatorname{adj}(Q)t$ の三成分であり、
三固定点における評価写像は同型である。
\end{corollary}

\begin{proof}
最初の等式と五次式の基底は\eqref{eq:moving-HB-v28}から従う。
標準形の第 $i$ 頂点で $t(e_i)=Q(e_i)e_i$ が成り立つため、
\[
 \operatorname{adj}(Q(e_i))t(e_i)=\det(Q(e_i))e_i.
\]
$C$ は三頂点を避けるので、これら三つの非零スカラーは全て可逆である。
従って五次式の評価行列は可逆な対角行列となり、後半も従う。
\end{proof}

\begin{corollary}[補助六次曲線の直接式]
\label{cor:bordered-gram-v28}
同じ標準形で $F=\det Q$ と置くと、補助Gram六次曲線は
\begin{equation}\label{eq:bordered-gram-v28}
 \det A\cdot\det(RA^{-1}R^t)
 =\det\begin{pmatrix}Q&t\\t^t&s\end{pmatrix}
 =sF-t^t\operatorname{adj}(Q)t
\end{equation}
で与えられる。
座標に依存しない式では $s=v^tAv$ と置く。
従って補助六次曲線も、$Q$ と標準的な直線 $\ell$ から直接得られる。
\end{corollary}

\begin{proof}
$M=[P^t\mid e_0]$ と置くと $RM=0$ である。
\eqref{eq:canonical-P-R-v28}の小行列式を取れば、$R$ の三次小行列式と
$M$ の相補的な四次小行列式が、相補添字の符号と一つの共通符号を除いて一致する。
Jacobiの相補小行列式恒等式とCauchy--Binetの公式を適用して
\[
 \det A\cdot\det(RA^{-1}R^t)=\det(M^tAM)
\]
を得る。右辺の最後の行・列による展開が\eqref{eq:bordered-gram-v28}である。
$v$ をスカラー倍すると右辺全体がそのスカラーの二乗倍となるだけなので、曲線は不変である。
\end{proof}

\paragraph{再現入口と適用範囲.}
\nolinkurl{intrinsic_sextic_normal_form_v28.py} は、$P,R$ の恒等式、
十九変数の線形解、九つの逆行列成分、基点における余因子Jacobi行列 $I_9$、
既存の整数基点との同一性、およびHilbert--Burch次数を厳密に検証する。
標準形の記述は新たな低精度近似を用いない。
この結果は七次元の選択した像に明示的な方程式と標準的な切断を与えるが、
$\mathcal Y_0$ の大域的有理性、九余因子軌跡の他の成分、
反対合の標識を除いたTorelli問題は別の問題として残る。

\section{標準五次netと七次平面被覆}
\label{sec:canonical-septic-cover-v29}

標準切断の十八点除数は、その支持を復元するだけでなく、
新しい平面への写像を定める。
本節では、前節と同じ選択成分 $\mathcal Y_0$ 上で、この写像の次数と
幾何学的monodromyを決定する。
十八単純基点を持つ五次netの交点数自体は通常の交点理論から従う。
ここでの内容は、netがPrym対から標準的に得られることと、
九余因子条件を満たす同一の成分の中で完全対称群 $\mathfrak S_7$ が生じることである。

\begin{theorem}[標準五次netによる次数七写像]
\label{thm:canonical-septic-cover-v29}
定理\ref{thm:canonical-moving-section-v28}の設定を用い、
$D$ は被約で、一次式列 $t$ の係数行列は可逆であるとする。
$F=\det Q$、$g=\operatorname{adj}(Q)t$ と置く。
このとき次が成り立つ。
\begin{enumerate}
 \item $g$ の三成分は $H^0(\mathbb P^2,\mathcal I_D(5))$ の基底であり、
 そのnetの基点スキームは正確に $D$ である。
 \item $\pi:X=\operatorname{Bl}_{D}\mathbb P^2\to\mathbb P^2$ とし、
 $H$ を直線の引き戻し、$E_1,\ldots,E_{18}$ を例外曲線とする。
 五次netは基点を持たない線束
 \[
  L=5H-\sum_{i=1}^{18}E_i
 \]
 を定め、全射かつ一般有限な次数七の射
 \begin{equation}\label{eq:septic-morphism-v29}
  f:X\longrightarrow
  \mathbb P\bigl(H^0(\mathbb P^2,\mathcal I_D(5))^*\bigr)
  \simeq\mathbb P^2
 \end{equation}
 を与える。
 \item $C$ の狭義変換を $\widetilde C$ とすると、
 $f|_{\widetilde C}$ は完全線形系 $|\eta(2H)|$ で与えられる。
 この線束の次数は十二、切断数は三である。
 \item $df$ の行列式で定まる分岐の源側の因子を $R_f$ とすると、
 \begin{equation}\label{eq:septic-ramification-class-v29}
  R_f\sim K_X+3L
      =12H-2\sum_{i=1}^{18}E_i
      =2\widetilde C,
  \qquad \deg(f_*R_f)=R_f\cdot L=24.
 \end{equation}
 後者は分岐の重複度を保った押し出しcycleの次数である。
\end{enumerate}
これらの仮定は $\mathcal Y_0$ の非空開部分で同時に満たされる。
復元開部分では、写像は抽象Prym対 $(C,\eta)$ から、源と標的の射影座標変換を除いて決まる。
\end{theorem}

\begin{proof}
余因子の恒等式より $Qg=Ft$ である。
$t$ の係数行列が可逆なので、射影平面の各点でその成分の一つは単元となる。
従って各局所環で $F\in(g_0,g_1,g_2)$ であり、
\eqref{eq:moving-ideal-v28}からイデアル層として
\[
 (g_0,g_1,g_2)^{\sim}
  =(F,g_0,g_1,g_2)^{\sim}=\mathcal I_D
\]
を得る。これは支持だけの等式ではない。
切断空間については\eqref{eq:moving-HB-v28}から次元三と上記の基底を得る。

$D$ は十八個の単純点なので、そのイデアルのblow-upでnetの基点は全部解消され、
引き戻したイデアルは $\mathcal O_X(-\sum E_i)$ となる。
従って $L$ は大域切断で生成される。
$L^2=25-18=7>0$ なので、対応する射の像は曲線にも点にもならない。
像は射影平面全体であり、$L=f^*\mathcal O(1)$ から一般次数は七である。

$\widetilde C\simeq C$ であり、標準切断の零除数は
$\mathcal O_C(D)=K_C\otimes\eta=\eta(3H)$ を満たす。
従って
\[
 L|_{\widetilde C}=\mathcal O_C(5H-D)=\eta(2H).
\]
制限の完全列
\[
 0\longrightarrow\mathcal O_{\mathbb P^2}(-1)
 \longrightarrow\mathcal I_D(5)
 \longrightarrow i_*\eta(2H)\longrightarrow0
\]
と $H^0(\mathcal O(-1))=H^1(\mathcal O(-1))=0$ により、
切断空間の制限は同型である。
特に得られる線形系は完全である。
$C\setminus D$ 上では $Q(x)g(x)=0$ なので、
この制限写像は $Q(x)$ の核直線を送る写像にも一致する。

標数零なので一般有限な射は可分であり、$\det(df)$ は恒等的に零ではない。
その零因子の線束は $\omega_X\otimes f^*\omega_{\mathbb P^2}^{-1}$ である。
$K_X=-3H+\sum E_i$ と $[\widetilde C]=6H-\sum E_i$ を代入すると
\eqref{eq:septic-ramification-class-v29}の線形同値を得る。
射影公式は押し出しcycleの次数を $60-36=24$ と与える。
収縮される成分はこの押し出しに寄与しない。

係数行列 $t$ は $A_0$ で単位行列であり、その可逆性は
既約成分 $\mathcal Y_0$ 上の非空開条件である。
前節で得た滑らかな $C$ と被約な $D$ の非空開条件と交差させられる。
最後に、標準切断の直線と $D$ はPrym対から復元され、
平面六次曲線の $H$ も抽象曲線から復元される。
$\mathbb P^2$ の十八点のblow-upと完全な五次netを用いる構成なので、
新たな座標や基底の選択は必要ない。
\end{proof}

\paragraph{二つの区別.}
式\eqref{eq:septic-morphism-v29}は一般有限な射である。
この節で扱う範囲には曲線を収縮する特殊化も含む。
一般パラメータでの有限平坦性は第\ref{sec:septic-finiteness-v30}節で証明する。
関数体の次数七拡大、あるいは有限となる標的開部分への制限が、以下で扱う被覆である。
また\eqref{eq:septic-ramification-class-v29}は線形同値であって
$R_f=2\widetilde C$ という因子の等式ではない。
被約分岐曲線そのものの次数二十四と既約性は、第\ref{sec:septic-branch-integrality-v30}節で証明する。

\subsection{一つの単純foldを持つ完全なファイバー}

\begin{lemma}[同じ成分上の厳密な特殊化]
\label{lem:septic-witness-v29}
次の六変数族は、可逆行列である範囲で $\mathcal Y_0$ に含まれる。
\[
 A=A_0+\operatorname{diag}(0,a,c,h,j,n,q),
 \qquad
 \det A=(ah-1)(cn-1)(jq-1).
\]
特に
\begin{equation}\label{eq:septic-Astar-v29}
 (a,c,h,j,n,q)=(3,-1,4,-2,5/4,3/4)
\end{equation}
で定める $A_*$ は同じ選択成分の滑らかな点である。
この点で $C$ は滑らかでスペクトル三角形の三頂点を避け、$D$ は十八個の相異なる点であり、$t=x$ である。
\end{lemma}

\begin{proof}
添字を $0;(1,3);(2,5);(4,6)$ の順に並べると、上の行列は
一つの一次ブロックと三つの二次ブロックに分かれる。
その逆行列の九つの指定成分は恒等的に零である。
可逆な範囲は六次元アフィン空間の非空開部分であり、既約で $A_0$ を含む。
$A_0$ では九方程式の零点集合が滑らかで、通る既約成分が一つしかないので、
この族全体は $\mathcal Y_0$ に含まれる。
これは二点でJacobi階数が同じであることから成分を推測する議論ではない。
$A_*$ では $\det A_*=495/8$ であり、九本の逆行列成分の
$(p,m,i,r,o,l,b,g,f)$ に関するJacobi小行列式は
$-4096/60037250625\ne0$ となる。

$x=(x_0,x_1,x_2)^t$、$x^{[2]}=(x_0^2,x_1^2,x_2^2)^t$ とし、
\begin{equation}\label{eq:septic-Mstar-v29}
 M=\begin{pmatrix}-1&3&-1\\4&-1&-2\\5/4&3/4&-1\end{pmatrix},
 \qquad (F_0,F_1,F_2)^t=Mx^{[2]}
\end{equation}
と置く。このとき
\begin{equation}\label{eq:septic-rank-one-form-v29}
 Q=\operatorname{diag}(F_0,F_1,F_2)+2xx^t,\qquad t=x,
 \qquad
 g=(x_0F_1F_2,\ x_1F_0F_2,\ x_2F_0F_1)^t.
\end{equation}
従ってこの特殊化でも $Q$ は対角行列ではない。
$\det M=-9/4$ であり、その全成分と二次小行列式は非零である。
後者は行対・列対を $(01),(02),(12)$ の順に並べると
\[
 \begin{pmatrix}
 -11&6&-7\\ -9/2&9/4&-9/4\\17/4&-3/2&5/2
 \end{pmatrix}
\]
となる。

式\eqref{eq:septic-rank-one-form-v29}の基点は、各対 $F_i=F_j=0$ の四点、
および各 $x_i=F_i=0$ の二点で尽くされる。
実際、三つの $F_i$ は同時には零にならず、一つだけが零なら対応する $x_i$ も零になる。
二次小行列式の非消滅は、前者の四点で全座標が非零であること、
異なる集合が交わらないこと、および二つの方程式の横断性を保証する。
後者の二点では残る二つの座標と二つの $F_j$ が非零であり、横断性も明らかである。
各局所基点イデアルは、それぞれ $(F_i,F_j)$ または $(x_i,F_i)$ である。
従って全基点スキームは被約で、長さは $3\cdot4+3\cdot2=18$ となる。

滑らかさには次の整数係数六次式を用いる。
\begin{align*}
 4\det Q={}&20x_0^6+57x_0^4x_1^2-78x_0^4x_2^2
 +42x_0^2x_1^4+15x_0^2x_1^2x_2^2+18x_0^2x_2^4\\
 &+9x_1^6-57x_1^4x_2^2-6x_1^2x_2^4+8x_2^6.
\end{align*}
三座標点での値は $20,9,8$ である。
各チャート $x_i=1$ において、この式の三偏微分が生成するイデアルの
$\mathbb F_{101}$ 上のGr\"obner基底は $(1)$ である。
これは添付の厳密検証で三チャートを独立に再計算する。
従って $101$ での還元は滑らかであり、標数零でも滑らかである。
\end{proof}

\begin{lemma}[一つの二重点と五つの単純点]
\label{lem:septic-simple-fold-v29}
$A_*$ に対応する射 $f$ の $[1:1:1]$ 上のファイバーは、
$[1:1:1]$ にある一つの単純foldと、五つの\'etale点からなる。
例外曲線上にも、無限遠にも、他の逆像はない。
\end{lemma}

\begin{proof}
まず $x_iF_i\ne0$ の開部分でファイバーを計算する。
$g=[1:1:1]$ は、ある非零スカラー $\lambda$ に対する
$x_i=\lambda F_i(x)$ と同値である。
代表元を $X=\lambda x$ と一意に正規化すると、三つのアフィン方程式
\begin{equation}\label{eq:septic-fixed-points-v29}
 \mathcal F(X):=X-MX^{[2]}=0
\end{equation}
を得る。これは対応する局所化されたファイバースキームの同型である。
斉次化の無限遠では $MX^{[2]}=0$ となり、$M$ が可逆なので解はない。
アフィン原点は射影点を与えず、そこでのJacobi行列は単位行列なので長さ一である。

辞書式順序 $X_0>X_1>X_2$ での厳密なGr\"obner基底は
\[
 X_0-u(z),\quad X_1-v(z),\quad
 \frac{z(z-1)^2q_5(z)}{6561},\qquad z=X_2,
\]
という形を持つ。ここで
\begin{equation}\label{eq:septic-qfive-v29}
 q_5(z)=6561z^5-115182z^4+650187z^3
       -1147860z^2-216212z-9948.
\end{equation}
七次式 $u,v$ の全係数は随伴certificateに記録し、元の三方程式から基底を再計算する。
必要な厳密な値は
\[
 \gcd(q_5,q_5')=1,\quad q_5(0)=-9948,\quad q_5(1)=-832454,
 \quad u(1)=v(1)=1
\]
および、$h=(z-1)^2q_5$ に対する
\begin{equation}\label{eq:septic-coordinate-gcd-v29}
 \gcd(u,h)=\gcd(v,h)=\gcd(z,h)=1
\end{equation}
である。従って原点を除くスキームの長さは七で、
$(1,1,1)$ に長さ二の点があり、他の五点は相異なる単純点である。
全ての座標と $F_i$ はそこで非零である。

次に、この計算から漏れる例外曲線がないことを確認する。
$F_i=F_j=0$ 型の基点では、第三成分 $g_k$ が二次以上で消えるので、
例外曲線の像は標的の座標直線 $y_k=0$ に含まれる。
$x_i=F_i=0$ 型の基点 $p$ では、局所パラメータ $(x_i,F_i)$ に関する一次部分から、
例外曲線の像は
\[
 p_kF_j(p)y_j-p_jF_k(p)y_k=0
\]
という直線になる。
この直線が $[1:1:1]$ を含む条件は
$p_j/F_j(p)=p_k/F_k(p)=\lambda$ である。
この場合、$X=\lambda p$ は\eqref{eq:septic-fixed-points-v29}の非零解で
$X_i=0$ を満たす。逆も同じ計算で成り立つ。
従って\eqref{eq:septic-coordinate-gcd-v29}がこの可能性を正確に排除する。
例外曲線以外で $g$ の三成分が非零なら、全ての $x_i,F_i$ も非零であり、
既に計算した開部分に含まれる。これでファイバーの完全性を得る。

残る二次条件を確認する。$(1,1,1)$ での $\mathcal F$ のJacobi行列は
\[
 J=\begin{pmatrix}3&-6&2\\-8&3&4\\-5/2&-3/2&3\end{pmatrix},
 \qquad\operatorname{rank}J=2.
\]
右核と左核の生成元を
\[
 v_0=(10/13,28/39,1)^t,\qquad w_0=(-1/2,-1/2,1)^t
\]
とすると、二次微分の組合せは
\begin{equation}\label{eq:septic-fold-Hessian-v29}
 w_0^t\,d^2\mathcal F(v_0,v_0)=-679/1521\ne0
\end{equation}
となる。
標的を $[u_0:u_1:1]$ とすると、近傍のファイバーは
$X=\operatorname{diag}(u_0,u_1,1)MX^{[2]}$ で与えられる。
その標的方向の微分は $-e_0,-e_1$ であり、左核との組合せは
$(1/2,1/2)$ である。従って全方程式は滑らかな曲面を定め、
標的への射の核方向における二次微分が非零である。
これは単純foldの局所形 $(s,t)\mapsto(s,t^2)$ の判定である。
階数二という一次の判定だけではなく、\eqref{eq:septic-fold-Hessian-v29}を用いている。
他の五点ではファイバースキームが単純なのでJacobi行列は可逆であり、射は\'etaleである。
\end{proof}

\subsection{幾何学的monodromy}

\begin{theorem}[完全対称群の実現]
\label{thm:septic-Sseven-v29}
$\mathcal Y_0$ のある非空Zariski開部分では、
定理\ref{thm:canonical-septic-cover-v29}の次数七拡大の幾何学的monodromy群は
$\mathfrak S_7$ である。
従って同じ結論は、七次元の選択されたPrym像の一般点で成り立つ。
\end{theorem}

\begin{proof}
まず $A_*$ で考える。
源 $X$ は既約なので、有限\'etaleとなる標的開部分の七枚のsheet上の作用は推移的である。
補題\ref{lem:septic-simple-fold-v29}のファイバーの周りでは、
二つのsheetだけが単純に入れ替わり、残る五つは\'etaleである。
従ってmonodromy群は一つの互換を含む。

次数七は素数なので、推移的な作用は原始的である。
ここで、原始的な置換群が互換を含めば完全対称群であることを確認する。
その互換の全共役を辺とする七頂点のグラフを作る。
連結成分は群によって置換されるblockをなす。
辺が存在し、作用が原始的なので、このグラフは連結である。
連結グラフの辺の互換は全対称群を生成するため、群は $\mathfrak S_7$ に等しい。
推移性と一つの互換だけから、任意の次数で完全対称群を結論しているわけではない。

最後に、このfoldが一般のパラメータにも残ることを説明する。
$A_*$ は $\mathcal Y_0$ の滑らかな点である。
その滑らかな開近傍 $B$ を縮めると、$t$ は可逆、$C$ は滑らか、
$D\to B$ は次数十八の有限\'etale族となる。
この相対的な $D$ に沿うblow-up $\mathcal X\to B$ は滑らかな曲面族であり、
五次netは $\mathcal X\to\mathbb P^2_B$ を定める。
上で計算したファイバーは零次元なので、固有性によりその標的近傍ではこの射は有限である。
源と標的は滑らかで次元が等しいから、そこでは有限平坦であり、階数は七である。

この有限平坦代数は、$[1:1:1]$ の標的\'etale近傍で、
階数二の一因子と階数一の五因子に分かれる。
階数一の因子は\'etaleである。
標数零の階数二因子は平方完成により $w^2=a$ と書け、
補題\ref{lem:septic-simple-fold-v29}の二次条件は
分岐式 $a$ の相対微分 $d_{/B}a$ が非零であることを与える。
これは開条件であり、$a=0$ はその近傍で $B$ 上滑らかに写る。
滑らかな射は開写像なので、$B$ のある非空開部分の各点で、
単純fold一個と\'etale点五個からなるファイバーが存在する。
各パラメータの源は依然として既約で次数は七であるから、上と同じ群論が適用できる。

得た開部分と既存のPrym復元開部分は、同じ既約成分内の非空開部分として交わる。
この性質は源・標的の座標変換で変わらないため、指定された商の一般点にも降りる。
\end{proof}

\paragraph{検証資料と未決問題.}
\nolinkurl{canonical_septic_cover_v29.py} と随伴JSONは、
同一成分を与える六変数族、九方程式のJacobi小行列式、
$Q$ と $g$ の恒等式、十八基点の非消滅条件、六次曲線の三チャートの滑らかさ、
完全な消去基底、例外曲線を排除する三つのgcd、
およびfoldの二次微分を厳密算術で再計算する。
通常実行は保存済みcertificateを変更せず照合する。
大域的な $\mathcal Y_0$ の有理性と全九余因子軌跡の成分分解は未決である。
一般の被約分岐曲線の次数と整性、分岐源の滑らかさは後の三節で決定する。
分岐曲線の特異点型の完全な分類は残る問題である。

\section{七次射の有限性と標準偏極の豊富性}
\label{sec:septic-finiteness-v30}

前節の七次射は、一般点で有限であることまでしか用いなかった。
ここでは、同じ選択成分 $\mathcal Y_0$ の非空開部分上で、
曲面全体にわたり有限となることを証明する。
この区別は実質的である。単純分岐を検出した六変数部分族には、
実際に三本の収縮曲線がある。

\begin{theorem}[標準七次射の有限平坦性]
\label{thm:septic-finiteness-v30}
選択成分 $\mathcal Y_0$ のある非空開部分に対し、標準七次射
\[
 f:X=\operatorname{Bl}_{D}\mathbb P^2\longrightarrow\mathbb P^2
\]
は有限平坦で次数七である。したがって
\[
 L=5H-\sum_{i=1}^{18}E_i=f^*\mathcal O_{\mathbb P^2}(1)
\]
は豊富である。この性質は源と標的の射影座標変換で不変なので、
指定された七次元のPrym像の非空開部分にも成立する。
\end{theorem}

\subsection{収縮曲線の次数に対する一様な制限}

以下、既存の良い開部分上で $D$ は相異なる十八点、
$C$ は滑らかな六次曲線、$\eta$ はその二位数線束とする。
$D\sim K_C\otimes\eta$ および
\[
 h^0(C,\eta(H))=0,\qquad
 h^0(C,\eta(2H))=3
\]
であり、後者の完全線形系は基点を持たない。

\begin{lemma}[四次式に対するCayley--Bacharach性]
\label{lem:septic-CB4-v30}
$D$ の任意の十七点を通る四次曲線は存在しない。
\end{lemma}
\begin{proof}
$p\in D$ とする。平面上の制限完全列より
\[
 H^0(\mathbb P^2,\mathcal I_{D-p}(4))
 \simeq H^0(C,\eta(H+p)).
\]
$\deg\eta(H+p)=7$、$g(C)=10$ である。
Riemann--Rochと $\eta(2H)$ の基点自由性から
\[
 h^0(C,\eta(H+p))
 =h^0(C,\eta(2H-p))-2=2-2=0.
\]
\end{proof}

\begin{lemma}[収縮候補の有限な型]
\label{lem:septic-contraction-types-v30}
$f$ が収縮する既約曲線は、次のいずれかの平面曲線の狭義変換である。
\begin{enumerate}
\item $D$ の五点を通る直線。
\item $D$ の十点を通る二次曲線。
\item $D$ の十五点を通る三次曲線。
\item $D$ の一点で重複度二、別の十三点で重複度一を持つ三次曲線。
\end{enumerate}
\end{lemma}
\begin{proof}
$L\cdot E_i=1$ なので、例外曲線は収縮されない。
平面曲線の次数を $d>0$、十八点での重複度を $m_i\ge0$ と書く。
収縮条件は $\sum m_i=5d$ である。
その狭義変換の算術種数は非負なので、随伴公式から
\[
 \sum m_i^2\le d^2+2d+2.
\]
Cauchyの不等式と合わせると
\[
 \frac{25d^2}{18}\le d^2+2d+2,
 \qquad 7d^2-36d-36\le0,
\]
従って $d\le6$ である。
$d=5$ のとき、和が二十五である十八個の非負整数の平方和の最小値は三十九、
$d=6$ のときは五十四であり、上の上界三十七、五十にそれぞれ反する。
$d=4$ のとき
\[
 \sum_i m_i(m_i-1)\le6
\]
である。重複度が正の点の個数を $n$ とすると
$2(20-n)\le6$、従って $n\ge17$ となり、
補題\ref{lem:septic-CB4-v30}に反する。
$d=1,2,3$ について同じ平方和の上界を用いると、列挙した四つの型だけが残る。
\end{proof}

\subsection{六変数部分族の収縮曲線}

既存の同一成分の六変数部分族で
\[
 M=\begin{pmatrix}-1&a&c\\h&-1&j\\n&q&-1\end{pmatrix},\qquad
 (F_0,F_1,F_2)^t=M(x_0^2,x_1^2,x_2^2)^t
\]
と置く。$Q=\operatorname{diag}(F_0,F_1,F_2)+2xx^t$、$t=x$ であり、
\[
 (g_0,g_1,g_2)=(x_0F_1F_2,x_1F_0F_2,x_2F_0F_1).
\]
$M$ の成分、二次小行列式、行列式がすべて非零となる開部分を考える。
十八基点は十二個の $F_i=F_j=0$ と、六個の $x_i=F_i=0$ からなり、
すべて横断的である。

\begin{lemma}[三本の $(-6)$ 曲線]
\label{lem:septic-three-conics-v30}
この六変数部分族で、$f$ が収縮する既約曲線は
$F_0=0,F_1=0,F_2=0$ の狭義変換 $T_0,T_1,T_2$ に限る。
これらは互いに交わらない滑らかな有理曲線であり、
\[
 T_i^2=-6,\qquad L\cdot T_i=0.
\]
各 $T_i$ の像は標的の第 $i$ 座標点である。
\end{lemma}
\begin{proof}
各 $F_i$ は相異なる十基点を通る滑らかな二次曲線なので、
自己交叉数と収縮の主張は直ちに従う。
$F_i$ と $F_j$ の四交点はすべて一度ずつblow-upされるため、狭義変換は交わらない。
例外曲線は $L$ との交点数一なので収縮されない。

標的の三座標が非零の点 $[u_0:u_1:u_2]$ を取る。
$\prod x_iF_i\ne0$ の部分でファイバーの方程式を射影的に正規化すると
\[
 Y=\operatorname{diag}(u_0,u_1,u_2)M Y^{[2]}
\]
となる。この二次方程式系を斉次化したとき、無限遠の方程式は
$\operatorname{diag}(u)M Y^{[2]}=0$ であり、$\det M\ne0$ により解を持たない。
従ってこのアフィン方程式系は零次元である。

標的の一つの座標だけが零なら、$F_i=0$ を除いた源の逆像は
対応する座標直線 $x_i=0$ 上にある。その直線上の写像は、残りの添字を $j,k$ として
\[
 [x_j:x_k]\longmapsto
 [x_j(m_{kj}x_j^2-x_k^2):x_k(-x_j^2+m_{jk}x_k^2)]
\]
であり、仮定した非消滅条件により三次の非定数写像である。
標的が座標頂点なら、上記の三二次曲線の外での逆像は源の座標頂点に限られる。
以上から他の曲線は収縮されない。
\end{proof}

\subsection{同じ成分内での三本の非有効化}

前版で用いた有理点
\[
 (a,c,h,j,n,q)=(3,-1,4,-2,5/4,3/4),\qquad s=1
\]
を $A_*$ と書く。他の正規形パラメータは零である。
従属変数を $(p,m,i,r,o,l,b,g,f)$ とする九逆行列条件のJacobi行列式は
\[
 -\frac{4096}{60037250625}\ne0.
\]
従って $A_*$ は同じ $\mathcal Y_0$ の滑らかな点であり、
自由変数 $d,e,k$ の微分をそれぞれ一、他の自由変数の微分を零とする接ベクトルが存在する。
その非零の従属変数の微分は次のとおりである。
\[
\begin{array}{c|l}
 \text{自由方向}&\text{従属変数の微分}\\\hline
 d&\dot o=15/16,\quad\dot l=5/4,\quad\dot f=3/4\\
 e&\dot p=4,\quad\dot i=-8,\quad\dot g=-2\\
 k&\dot m=-1,\quad\dot r=3,\quad\dot b=-3
\end{array}
\]
これらは $\dot Z=-Z\dot A Z$ を九つの禁止成分で零とすることで計算される。

\begin{lemma}[三つの二次曲線条件への非接性]
\label{lem:septic-conic-obstruction-v30}
$A_*$ の近傍で十八基点を\'etale局所的に標識付ける。
$F_i=0$ 上にある十点が引き続き二次曲線上にあるという条件は、
$i=0,1,2$ のいずれについても $\mathcal Y_0$ の真の閉条件である。
\end{lemma}
\begin{proof}
十点の二次式評価行列を $E_i$ とする。$A_*$ では
$\ker E_i=\langle F_i\rangle$ である。
接ベクトルに沿う基点の速度を $\dot p_\nu$ とすれば、
二次式 $F_i+\epsilon H$ が動く十点を通るための必要十分な一次条件は
\[
 E_i H=b_i,\qquad
 (b_i)_\nu=-\nabla F_i(p_\nu)\cdot\dot p_\nu.
\]
基点の速度は $dg(p_\nu)\dot p_\nu+\dot g(p_\nu)=0$ を
射影アフィンチャートで解いて求める。

この線形系を $\mathbb F_{7561}$ 上で計算すると、十八基点はすべて分裂し、
各 $E_i$ の階数は五である。一方、付加行列の階数は
\[
\begin{array}{c|ccc}
 &d&e&k\\\hline
 F_0&6&6&5\\
 F_1&6&5&6\\
 F_2&5&6&6
\end{array}
\]
となる。従って $F_0,F_1$ には $d$ 方向、$F_2$ には $e$ 方向の一次延長が存在しない。

有限体計算から標数零への移行を明記する。
$A_*$、三接ベクトル、逆関数定理に使うJacobi小行列式の分母はすべて
$7561$ において単元である。十八基点はこの素数においても相異なり、
各基点で二つの局所方程式のJacobi行列は階数二である。
従って基点スキームは対応する局所整数環上で有限\'etaleである。
各十点への二次式評価写像も整数環上で定義され、
その特殊ファイバーは階数五、核は $F_i$ の直線である。
$5\times5$ 小行列式の一つが単元なので、二次式の五係数を五点の条件で消去できる。
残りの条件は底の正則関数となり、上の階数六はその微分の少なくとも一つが非零であることを示す。
これは標数零でも非零であり、当該閉条件は成分全体ではない。

\nolinkurl{septic_finiteness_v30.py} は、接ベクトル、十八点の座標、
速度、三つの評価行列、付加行列の非零小行列式を厳密に再計算する。
随伴JSONには行・列の添字と小行列式の値も保存する。
\end{proof}

\begin{proof}[定理\ref{thm:septic-finiteness-v30}の証明]
十八点を標識付ける有限\'etale被覆のうち、$A_*$ の上で選んだ点を通る既約成分 $S$ 上で考える。
補題\ref{lem:septic-contraction-types-v30}に現れる重複度型は有限個である。
各型について、指定された点でその重複度以上に消える次数 $d\le3$ の曲線が
存在するという条件は、平面曲線の射影係数空間との閉incidenceの固有像である。
従って閉条件である。

$A_*$ でそのような曲線が存在するなら、その全狭義変換と $L$ の交点数は零以下である。
$L$ はnefなので実際には零であり、全既約成分が収縮される。
補題\ref{lem:septic-three-conics-v30}から、それらは三つの二次曲線に限られる。
次数は三以下なので、候補は各 $F_i$ 自身のみである。
従って、$A_*$ を含む開近傍では、閉incidenceとして残る可能性はこの三つの十点二次曲線条件だけである。
補題\ref{lem:septic-conic-obstruction-v30}により、三条件はいずれも真の閉条件である。
既約な $S$ からこれらを除く非空開部分では、収縮曲線は存在しない。
この性質は標識付けに依存せず、\'etale写像は開写像なので、その像は $\mathcal Y_0$ の非空開部分を与える。

固有な全射 $f$ は曲線ファイバーを持たないので準有限、従って有限である。
滑らかな曲面から滑らかな曲面への有限射は、Cohen--Macaulay性と次元公式によって平坦である。
その次数は既に計算した $L^2=7$ である。
有限射による豊富な線束の引戻しは豊富なので、$L$ の豊富性も従う。
\end{proof}

\paragraph{この段階での分岐の数値.}
有限平坦な一般射でも、前節の
\[
 R_f\sim2\widetilde C,\qquad R_f^2=72,\qquad K_X\cdot R_f=0
\]
はそのまま成り立ち、分岐押出しサイクルの次数は二十四である。
分岐曲線の被約性と既約性は次節で証明し、分岐源の滑らかさは
定理\ref{thm:septic-ramification-smoothness-v30}で確定する。

\section{分岐曲線の既約性と次数二十四}
\label{sec:septic-branch-integrality-v30}

有限性の証明に使った特殊曲面には三本の収縮二次曲線が残る。
その退化を用いると、一般の分岐曲線の既約性も決定できる。
要点は、特殊な分岐源の中で $L$ と正に交わる成分が一本しかないことである。

\begin{theorem}[一般の分岐曲線]
\label{thm:septic-branch-integrality-v30}
同じ選択成分 $\mathcal Y_0$ のある非空開部分で、次が成り立つ。
\begin{enumerate}
\item 標準七次射 $f:X\to\mathbb P^2$ は有限平坦であり、分岐源 $R_f$ は幾何学的に整である。
\item 被約分岐曲線 $B_f\subset\mathbb P^2$ は幾何学的に整な二十四次曲線である。
写像 $R_f\to B_f$ は双有理であり、一般分岐点での慣性は互換である。
\item $p_a(R_f)=37$ であり、
\[
 \delta(B_f)=216+\delta(R_f)
\]
が成り立つ。ただし $\delta$ は正規化による算術種数の減少を表す。
次節の定理\ref{thm:septic-ramification-smoothness-v30}では、さらに
$R_f$ が滑らかであることを証明し、正規化種数三十七と $\delta(B_f)=216$ を得る。
\end{enumerate}
\end{theorem}

\subsection{既約な残余六次曲線を持つ特殊点}

六変数部分族で別の有理点
\[
 (a,c,h,j,n,q)=(1,2,3,4,5,6),\qquad s=1
\]
を選び、これを $A_\#$ と書く。この点は $A_0$ を通る既約な六変数族にあるため、
前と同じ選択成分 $\mathcal Y_0$ に属する。
ここでは
\[
 M_\#=\begin{pmatrix}-1&1&2\\3&-1&4\\5&6&-1\end{pmatrix},\qquad
 \det M_\#=92.
\]
$F=M_\#x^{[2]}$ とし、残余critical sexticを
\[
 J(x)=\det\bigl(\operatorname{diag}(F)-2\operatorname{diag}(x)M_\#\operatorname{diag}(x)\bigr)
\]
と置く。明示式は
\begin{align*}
 J={}&15x_0^6-22x_0^4x_1^2-67x_0^4x_2^2-43x_0^2x_1^4
      -502x_0^2x_1^2x_2^2-110x_0^2x_2^4\\
     &+6x_1^6-59x_1^4x_2^2-138x_1^2x_2^4+8x_2^6.
\end{align*}

\begin{lemma}[分岐源の被約な特殊ファイバー]
\label{lem:septic-special-ramification-v30}
$A_\#$ で $C=V(\det Q)$ と $V(J)$ はともに滑らかである。
$V(J)$ の狭義変換を $R_0$ とすると
\[
 R_{f_\#}=R_0+T_0+T_1+T_2
\]
はすべての成分の係数が一の被約除数である。
$R_0$ は滑らかな種数十の曲線であり、
\[
 [R_0]=6H-\sum_{p\in D_{\rm coord}}E_p,\qquad
 L\cdot R_0=24,\qquad L\cdot T_i=0,
\]
ただし $D_{\rm coord}=\bigcup_iV(x_i,F_i)$ は六点からなる。
\end{lemma}
\begin{proof}
三つの二次曲線と十八基点の横断性は、$M_\#$ の全成分と二次小行列式の非消滅から従う。
$C$ と $J$ の滑らかさは $\mathbb F_{101}$ 上で証明される。
両曲線について三つのアフィンチャートすべてで、三つの斉次偏微分からなるイデアルの
Gr\"obner基底が $\{1\}$ となる。
これは射影的な特異点がないことを示し、標数零での滑らかさも従う。
計算は \nolinkurl{septic_ramification_v30.py} と随伴JSONに記録する。

既存の恒等式
\[
 \det(dg)=-5F_0F_1F_2J
\]
を使う。$J$ は十二個の $F_i=F_j=0$ を避け、
六個の $x_i=F_i=0$ を通る。
前者は $F_i,F_j$ による交点の二乗座標を代入して非零を確かめ、
後者は $J|_{x_i=0}$ が $F_i|_{x_i=0}$ で割り切れることから従う。
滑らかさにより後者の六点での重複度は一である。
従って、十二点では $F_0F_1F_2J$ の重複度は二、残る六点でも二である。

単純基点の一度のblow-upで、平面上のJacobian除数の狭義変換に対する例外係数は
基点重複度から二を引いた値となる。
これは $g_i=uh_i$ と書くblow-upチャートで
$dg_0\wedge dg_1\wedge dg_2$ を計算することによっても得られる。
ここでは例外成分は残らず、主張した四成分の和がちょうど $R_{f_\#}$ である。
$J$ は滑らかな平面六次曲線なので、その狭義変換も滑らかで種数は十である。
交点数は
\[
 L\cdot R_0=5\cdot6-6=24
\]
である。
\end{proof}

\begin{proof}[定理\ref{thm:septic-branch-integrality-v30}の証明]
十八基点が被約である良い開部分上で相対blow-upを行い、
滑らかな曲面族 $\mathcal X$ と射 $\mathcal f$ を考える。
相対微分の行列式は各ファイバーで零でない切断なので、
分岐源 $\mathcal R$ は相対有効Cartier除数であり、底上平坦である。

幾何学的一般ファイバーにおいて、分岐源が複数成分を持つか、
ある成分の係数が二以上であると仮定する。
必要なら底の有限拡大を行って成分を定義し、
$A_\#$ への離散付値環を通じてそれらの有効サイクルを特殊化する。
特殊化は有効性、$L$ との交点数、和としての $\mathcal R$ を保つ。
一般ファイバーでは定理\ref{thm:septic-finiteness-v30}により $L$ は豊富なので、
各成分は $L$ と正に交わる。

一方、補題\ref{lem:septic-special-ramification-v30}の特殊ファイバーでは、
$L$ と正に交わる成分は $R_0$ のみである。
従って各一般成分の特殊化は $R_0$ を含まなければならない。
しかし特殊ファイバーにおける $R_0$ の係数は一である。
このため一般成分は一本、その係数も一に限られる。
幾何学的一般分岐源は整であり、平坦な固有族における幾何学的整性の開性により、
底を縮めれば全ファイバーでこの性質が成り立つ。

前節で証明した単純foldの持続開部分と交わらせる。
そのfoldの標的近傍で、有限平坦代数は分岐する階数二の因子と
\'etaleな五つの階数一因子に分かれる。
従ってその近傍の分岐曲線は滑らかであり、分岐源からはただ一つの点が一対一に写る。
既約な $R_f$ からその像 $B_f$ への有限写像の一般次数は、これにより一である。
$R_f$ は係数一の分岐源なので、一般分岐指数は二である。
既存の押出し計算
\[
 \deg f_*R_f=L\cdot R_f=24
\]
は、ここで被約な既約分岐曲線そのものの次数二十四を与える。

最後に $R_f\sim12H-2\sum E_i$ から
\[
 R_f^2=72,\qquad K_XR_f=0,\qquad p_a(R_f)=1+72/2=37.
\]
$B_f$ は二十四次なので $p_a(B_f)=23\cdot22/2=253$ である。
両曲線は双有理で同じ正規化を持つため
\[
 \delta(B_f)-\delta(R_f)=253-37=216
\]
となる。
\end{proof}

\begin{corollary}[六変数収縮部分族の分岐]
\label{cor:septic-contracted-branch-v30}
六変数部分族の一般点では、残余分岐源は滑らかな種数十の曲線である。
その像は既約な二十四次分岐曲線であり、残余分岐源からの写像はその正規化である。
\end{corollary}
\begin{proof}
補題\ref{lem:septic-special-ramification-v30}により、
残余六次曲線が滑らかである条件はこの既約な六変数族の非空開条件である。
前節の単純foldの持続論法は、滑らかな六変数部分族そのものを底にしても適用できる。
従って、同じ族の別の非空開部分で単純foldと五つの\'etale点からなる完全なファイバーが存在する。
二つの開部分は交わる。
唯一の正次数の分岐源は残余六次曲線なので、同じ次数一の議論により結論を得る。
\end{proof}

\paragraph{次節での強化.}
定理\ref{thm:septic-ramification-smoothness-v30}では、同じ選択成分内で
$R_f$ が滑らかであることを証明する。
これにより $R_f$ 自身が種数三十七の正規化となり、$\delta(B_f)=216$ が確定する。
さらに正規化写像の微分零因子の長さと双対曲線の次数を求める。
通常の節点・尖点だけが生じるという特異点型の限定は、これらとは別に残る問題である。

\section{分岐源の滑らかさと正規化種数三十七}
\label{sec:septic-ramification-smoothness-v30}

分岐源の整性に続いて、その滑らかさも同じ制約付き成分上で証明できる。
特殊ファイバーの三十交点を用いるが、六つの交点では一次変形だけでは判定できない。
これらは十八点のblow-upの例外曲線上にあり、基点を動かした上で二次まで調べる必要がある。

\begin{theorem}[分岐曲線の正規化]
\label{thm:septic-ramification-smoothness-v30}
選択成分 $\mathcal Y_0$ のある非空開部分で、標準七次射の分岐源 $R_f$ は
滑らかな幾何学的整曲線で、種数は三十七である。
従って $R_f\to B_f$ は二十四次分岐曲線の正規化であり、
\[
 g(\widetilde B_f)=37,\qquad \delta(B_f)=216.
\]
\end{theorem}

\subsection{特殊分岐源の三十個の節点}

前節の $A_\#$ と $R_{f_\#}=R_0+T_0+T_1+T_2$ を用いる。
この点における九逆行列条件の、従属変数
$(p,m,i,r,o,l,b,g,f)$ に関するJacobi行列式は
\[
 -\frac{1}{29376588816}\ne0
\]
である。従って $A_\#$ も $\mathcal Y_0$ の滑らかな点である。

\begin{lemma}[節点とその局所平滑化]
\label{lem:septic-thirty-node-smoothing-v30}
$R_{f_\#}$ の特異点はちょうど三十個の通常の節点である。
そのうち二十四個は平面上の基点の外にあり、六個は
$x_i=F_i=0$ の各基点の例外曲線上にある。
各節点の局所平滑化関数は $\mathcal Y_0$ 上で恒等的に零ではない。
\end{lemma}

\begin{proof}
以下の有限計算はすべて $\mathbb F_{163}$ 上で行い、最後に標数零への移行を説明する。
$A_\#$、逆関数定理のJacobi行列式、基点の局所方程式の必要な小行列式は
この素数で単元である。
前節の六次曲線 $C,J$ は、この素数でも三つのアフィンチャートの全てで滑らかである。

$T_i$ 同士は交わらず、$R_0\cdot T_i=10$ である。
平面の基点を除いた $F_i=J=0$ において、二つの勾配が従属となるイデアルを
$x_iF_jF_k$ で飽和すると、全チャートで単位イデアルとなる。
したがってこの部分の交点は全て横断的である。
$J$ は十二個の $F_i=F_j=0$ を避ける。
残る六個の基点では $J$ と $F_i$ の接線が一致するが、
一度のblow-up後には以下の二次式のHessianが非零となり、通常の節点を与える。
従って各 $T_i$ 上に基点外の八点と例外曲線上の二点があり、合計三十点となる。

まず、基点外の二十四点について一次の平滑化を確認する。
自由変数 $d,e,k$ を一つずつ動かし、他の自由変数を固定した接方向は
\[
\begin{array}{c|l}
 d&\dot o=30,\quad\dot l=5,\quad\dot f=6\\
 e&\dot p=3,\quad\dot i=12,\quad\dot g=4\\
 k&\dot m=2,\quad\dot r=1,\quad\dot b=2
\end{array}
\]
で与えられる。
$J_{\rm tot}=\det(dg)$ と置くと、その一次変化は
\[
 \dot J_{\rm tot}
   =\operatorname{tr}\bigl(\operatorname{adj}(dg)\,d\dot g\bigr)
\]
である。三つの方向に対するこれらの多項式を $\dot J_d,\dot J_e,\dot J_k$ と書く。
全ての $i$ と全てのアフィンチャートで
\[
 (F_i,J,\dot J_d,\dot J_e,\dot J_k):(x_iF_jF_k)^\infty=(1)
\]
となる。したがって基点外の各交点では、少なくとも一方向の平滑化係数が一次で非零である。

次に例外曲線上の六点を扱う。
ここでは変形の規約を
\[
 A(\epsilon)=A_\#+\epsilon D+\epsilon^2 E+O(\epsilon^3)
\]
とする。従って $E$ は二階微分の半分である。
$Z=A_\#^{-1}$ と置くと、逆行列の二次係数は
\[
 ZDZDZ-ZEZ
\]
である。九つの禁止成分を零にする条件を解くと、
上の単独の $d,e,k$ 方向はいずれも $E=0$ を持つ。
高次係数が零であるとは主張しない。
Jacobi行列式が単元なので、これらの二次jetは実際の形式的な曲線
$\operatorname{Spec}k[[\epsilon]]\to\mathcal Y_0$ に持ち上がる。

$i$ 以外の添字を $j<k$ とし、$x_j=1$ のチャートを使う。
動く基点を
\[
 p(\epsilon)=p+\epsilon p_1+\epsilon^2p_2+O(\epsilon^3)
\]
と書き、$g(\epsilon)=g+\epsilon g_1+\epsilon^2g_2+O(\epsilon^3)$ とする。
その係数は
\begin{align*}
 dg(p)p_1+g_1(p)&=0,\\
 dg(p)p_2+\frac12d^2g(p)[p_1,p_1]+dg_1(p)p_1+g_2(p)&=0
\end{align*}
から求まる。blow-upチャートを
\[
 x_i=p_i(\epsilon)+u,\qquad x_k=p_k(\epsilon)+uw,\qquad x_j=1
\]
と取る。このチャートで $g$ は $u$ で割り切れ、$h=g/u$ は解消された写像を表す。
問題の点では $h_i$ が単元である。従って局所分岐源の方程式として
\[
 r(\epsilon,u,w)=\det\begin{pmatrix}
 h_0&h_1&h_2\\
 \partial_u h_0&\partial_u h_1&\partial_u h_2\\
 \partial_w h_0&\partial_w h_1&\partial_w h_2
 \end{pmatrix}
\]
を用いてよい。

これを二次まで展開すると
\[
 r=\frac12\xi^tH\xi+\epsilon c^t\xi+b\epsilon^2+
       \text{三次以上の項},\qquad \xi=(u,w)^t
\]
となる。六点とも $H$ は可逆で、$\epsilon$ の一次定数項は零である。
$\partial_ur=\partial_wr=0$ を形式的に解くと
$\xi=-\epsilon H^{-1}c+O(\epsilon^2)$ なので、局所平滑化関数の二次係数は
\[
 b-\frac12c^tH^{-1}c
\]
である。その厳密な値は次のとおりである。座標と数値は $\mathbb F_{163}$ で読む。
\[
\begin{array}{c|c|c|c|c}
 \text{成分}&\text{基点}&\text{自由方向}&\det H&
 b-\tfrac12c^tH^{-1}c\\\hline
 T_0&(0,1,\pm9)&d&45&118\\
 T_1&(1,0,\pm23)&d&19&14\\
 T_2&(1,\pm61,0)&e&157&31
\end{array}
\]
従って六つの節点も、それぞれ同じ成分上の二次変形により平滑化される。

ここで有限体の計算と標数零の関係を確認する。
$R_0\cap(T_0\cup T_1\cup T_2)$ は、$163$ における局所整数環上で
次数三十の有限\'etaleスキームとなる。
例外曲線上の六点は、その閉かつ開な有限\'etale部分をなす。
従って残る二十四点も有限\'etaleであり、有限体への還元時に基点へ逃げることはない。
有限\'etale被覆を取って各点を標識付ければ、上で計算した平滑化係数は
局所整数環上の正則なjet係数である。
有限体で非零となる係数は標数零でも非零なので、
すべての節点の局所平滑化関数が成分上恒等零ではないことを得る。

\nolinkurl{septic_ramification_smoothness_v30.py} と随伴JSONは、
滑らかさ、全ての飽和イデアル、基点の一次・二次速度、六点のHessianと二次係数を再計算する。
六点での計算はさらに、Euler公式による
\[
 r=\frac{\operatorname{sgn}(j,i,k)}{5u^2}\,J_{\rm tot}
\]
を用いた独立な展開と一致することを確認する。
\end{proof}

\begin{proof}[定理\ref{thm:septic-ramification-smoothness-v30}の証明]
$A_\#$ の近傍で相対分岐源を考える。
その相対特異点スキームは固有であり、$A_\#$ 上のファイバーは上記の三十点だけである。
底を縮めれば相対特異点スキームは有限となる。
各節点の近傍で、源の二変数に関する微分方程式をHessianの可逆性によって解くと、
特異なファイバーの像は一つの局所平滑化関数の零点で与えられる。
補題\ref{lem:septic-thirty-node-smoothing-v30}により、三十点の全てでその関数は恒等零ではない。
従って相対特異点スキームに底を支配する成分はなく、
底の非空開部分で $R_f$ は滑らかとなる。

前節で得た幾何学的整性と合わせると、$R_f$ は滑らかな整曲線である。
随伴公式による $p_a(R_f)=37$ がその種数に等しい。
有限双有理写像 $R_f\to B_f$ は正規化であり、
$p_a(B_f)=253$ から $\delta(B_f)=253-37=216$ を得る。
\end{proof}

\begin{corollary}[正規化写像の微分零因子]
\label{cor:septic-nonimmersion-divisor-v30}
一般の $f$ について、正規化写像 $\nu:R_f\to B_f\subset\mathbb P^2$ の
微分が零となる有効除数の長さは七十二である。
これは重複度を込めた数であり、相異なる通常尖点の個数を主張するものではない。
\end{corollary}

\begin{proof}
$R=R_f$ と略記する。$df$ がある点で階数零なら、
その二次行列式の微分も零となり、分岐源が特異になる。
従って $R$ の全ての点で $df$ の階数は正確に一である。
$R$ は被約なので、$i:R\hookrightarrow X$ として
\[
 0\longrightarrow T_X\xrightarrow{df}f^*T_{\mathbb P^2}
 \longrightarrow i_*\mathcal M\longrightarrow0
\]
があり、$\mathcal M$ は $R$ 上の線束である。
曲面のRiemann--Rochを使う。
$\chi(\mathcal O_X)=1$、$K_X^2=-9$、$c_2(X)=21$、$L^2=7$、$K_XL=3$ より
\begin{align*}
 \chi(T_X)&=2+K_X^2-c_2(X)=-28,\\
 \chi(f^*T_{\mathbb P^2})
   &=2+\frac12(3L)(3L-K_X)-3L^2=8.
\end{align*}
従って $\chi(\mathcal M)=36$ である。
$g(R)=37$ なので $\deg\mathcal M=72$ となる。

$df|_R$ の像線束を $\mathcal I$ と書けば、
\[
 0\longrightarrow\mathcal I\longrightarrow f^*T_{\mathbb P^2}|_R
       \longrightarrow\mathcal M\longrightarrow0
\]
より $\deg\mathcal I=3L\cdot R-72=0$ である。
$d\nu$ は非零の線束準同型
\[
 T_R\longrightarrow\mathcal I
\]
を与える。非零性は標数零と $\nu$ の双有理性から従う。
その零除数の次数は
$\deg\mathcal I-\deg T_R=0-(2-2\cdot37)=72$ である。
\end{proof}

\begin{corollary}[双対分岐曲線]
\label{cor:septic-dual-branch-v30}
一般の分岐曲線 $B_f$ の双対曲線 $B_f^*\subset(\mathbb P^2)^*$ は
次数四十八であり、その正規化も種数三十七である。
\end{corollary}

\begin{proof}
$N=\nu^*\mathcal O_{\mathbb P^2}(1)$ とし、前系の微分零因子を $Z_\nu$ と書く。
$\deg N=24$、$\deg Z_\nu=72$ である。
$\nu$ を三つの局所切断 $(a_0,a_1,a_2)$ で表すと、接線の座標は
$a_i\,da_j-a_j\,da_i$ で与えられる。
これらは $K_R\otimes N^{\otimes2}$ の切断であり、共通零因子は正確に $Z_\nu$ である。
従ってGauss写像 $\gamma:R\to B_f^*$ は
\[
 \gamma^*\mathcal O_{(\mathbb P^2)^*}(1)
   \simeq K_R\otimes N^{\otimes2}(-Z_\nu)
\]
を満たし、その次数は $72+48-72=48$ となる。

標数零では $\gamma$ は双有理であることも確認できる。
一般点での接線を $\ell$ とすると、$\ell(\nu)=0$、$\ell(d\nu)=0$ から
$(d\ell)(\nu)=0$ が従う。
$\gamma$ は非定数かつ分離的なので、双対曲線の一般点 $\ell$ における接線は
元の点 $\nu$ を復元する。従ってGauss写像の一般次数は一である。
上の次数計算が $\deg B_f^*=48$ を与え、正規化種数も従う。
\end{proof}

\paragraph{分岐特異点の型はなお別の問題である.}
本節で確定した数値は、正規化種数三十七、$\delta$ の総和二百十六、
正規化写像の微分零因子の長さ七十二、および双対曲線の次数四十八である。
分岐曲線の特異点が通常の節点と尖点だけであることは、これとは別の主張である。
もしその性質が証明されれば、通常尖点が微分零因子に重複度一で寄与するため
尖点は七十二個、$\delta$ の残りから節点は百四十四個となる。
現時点ではこの特異点型の限定を仮定しない。

% 挿入位置: v25 の補助Gram曲線との二重交差命題の後。
% この節の計算は同一Hensel分枝を使い、点配置のmonodromyを主張しない。
\section{補助Gram曲線の通常二重点と既約性}
\label{sec:complementary-nodes-v26}
\label{subsec:complementary-gram-nodes-v26}

補助Gram曲線の固定点特異点は、全incidence方程式を用いると
通常二重点まで決定できる。この局所的結論には、新しい一般性仮定や計算を要しない。

\begin{proposition}[固定点における非退化な法線形式]
\label{prop:fixed-gram-node-v26}
第\ref{sec:prym-setup-v26}節の全incidence、可逆対称pairing、
平面全体での $\operatorname{rank}P=p$、有限被約な $Z$ の条件を満たすとする。
各固定点 $x_0$ において、$L(x_0)$ の行空間上の $A$ による対称形式は非退化である。
従って $P(x_0)AP(x_0)^t$ の階数は $p$、
$B(x_0)AB(x_0)^t$ の階数は $q-1$ であり、
$x_0$ は $\Gamma_B$ の通常二重点である。
\end{proposition}

\begin{proof}
平面の座標を変え、$x_0=[1:0:0]$ とする。
以下では $M=L_0=L(x_0)$、$L(1,s,t)=M+sL_1+tL_2$ と書く。
被約な孤立退化点では $\operatorname{rank}M=5$ である。
実際、階数が四以下なら六次小行列式の全ての一次微分が零となり、
局所的な退化イデアルは極大イデアルを生成できない。
左核の生成元を行ベクトル $u=(0,b)$ と取る。
固定点なので $uJ=-u$ であり、$uM=0$ である。

$w_i=uL_i=bB_i$ $(i=1,2)$ と置く。
全incidence方程式を左から $u$ で掛けると
\[
 uL_iAM^tJ=uJMAL_i^t=0
\]
となるので、
\begin{equation}\label{eq:fixed-normal-orthogonality-v26}
 w_iAM^t=0,\qquad Aw_i^t\in K:=\ker M
\end{equation}
を得る。ここで $\dim K=2$ である。

階数五の行列の退化locusに対する法線微分は
\[
 T_{x_0}\mathbb P^2\longrightarrow K^*,\qquad
 (a_1,a_2)\longmapsto (a_1w_1+a_2w_2)|_K
\]
である。局所的に五次可逆小行列を消去すれば、退化イデアルは残る一行二列の
二成分で生成されるので、$Z$ が被約であることから、この微分は同型となる。
従って $w_1,w_2$ は独立であり、式\eqref{eq:fixed-normal-orthogonality-v26}により
$Aw_1^t,Aw_2^t$ は $K$ の基底をなす。
この基底に関する法線微分の行列は
\begin{equation}\label{eq:fixed-normal-gram-v26}
 H_0=(w_iAw_j^t)_{1\leq i,j\leq2}
\end{equation}
なので、$H_0$ は非退化である。

行空間 $R=\operatorname{Row}M$ の直交補空間は
$\langle w_1,w_2\rangle$ であり、そこに制限した形式が
式\eqref{eq:fixed-normal-gram-v26}である。
従って $R$ も非退化である。
さらに
\[
 R=\operatorname{Row}P(x_0)\mathbin{\perp}\operatorname{Row}B(x_0)
\]
は、次元 $p$ と $q-1$ の直交直和であるから、両項の制限形式も非退化となる。

$B$ の行の定数基底変換により、最後の行を $bB$ とし、
他の $q-1$ 行を $C(s,t)$ とする。
$C(0,0)AC(0,0)^t$ は可逆であり、最後の行は $sw_1+tw_2$ である。
式\eqref{eq:fixed-normal-orthogonality-v26}により、
Gram行列の混合ブロックは $(s,t)$ の二次以上の項からなる。
Schur補行列を取ると、$\det(BAB^t)$ の二次初項は
\[
 \det\bigl(C(0,0)AC(0,0)^t\bigr)
 (s,t)H_0(s,t)^t
\]
となる。これは非退化な二次形式であるから、$x_0$ は通常二重点である。
\end{proof}

\begin{lemma}[一つの次数による射影的な長さの上界]
\label{lem:jacobian-length-bound-v26}
$k$ を無限体、$S=k[x,y,z]$、$I\subset S$ を斉次イデアルとする。
ある $D$ と $f\leq D$ に対し $\dim_k(S/I)_D\leq f$ なら、
$\operatorname{Proj}(S/I)$ は零次元または空であり、その長さは $f$ 以下である。
\end{lemma}

\begin{proof}
$I$ を飽和しても指定次数の商空間の次元は増えない。
$R=S/I^{\mathrm{sat}}$ とし、$R$ が零でなければ、
飽和性と $k$ の無限性により $R$ 上の非零因子となる一次式 $\ell$ を取れる。
$h_d=\dim_k R_d$、$h_{-1}=0$ と置くと
\[
 h_d-h_{d-1}=\dim_k(R/\ell R)_d\geq0.
\]
$R/\ell R$ は次数一で生成されるので、そのある次数の成分が零なら、
以後の成分も全て零である。
もし次数零から $D$ までが全て非零なら $h_D\geq D+1$ となり、
$h_D\leq f\leq D$ に反する。
従って $h_d$ は $D$ 以前に定常となり、その定常値は $f$ 以下である。
これは射影スキームが零次元または空で、その長さが $f$ 以下であることを意味する。
\end{proof}

\begin{theorem}[補助曲線の特異点の完全な決定]
\label{thm:complementary-gram-nodal-v26}
固定成分 $Z_1,Z_3$ は、定理\ref{thm:gram-smooth-v25}の条件に加えて
次の性質を持つ非空開部分を含む。
補助Gram曲線 $\Gamma_B$ は幾何学的に既約であり、
その特異点は指定反対合の固定点ちょうど全部で、全て通常二重点である。
\[
\begin{array}{c|c|c|c|c}
 \text{成分}&\deg C&\deg\Gamma_B&
 \Gamma_B\text{ の通常二重点数}&\text{正規化の種数}\\\hline
 Z_1&6&6&3&7\\
 Z_3&4&8&5&16
\end{array}
\]
非固定点では $C$ と $\Gamma_B$ はともに滑らかであり、
各点の局所交差数は二である。
従って $Z_1$ では十八点、$Z_3$ では十六点で通常の二重接触を持つ。
\end{theorem}

\begin{proof}
定理\ref{thm:gram-smooth-v25}と同じ代数的 $101$ 進Hensel点を使う。
基点、pivot変数、自由変数、seedは、被約二十一点配置の既存証明書と全て照合する。
$G=\det(BAB^t)$、$d=2q$、$D=3d-5$ とし、
$G_x,G_y,G_z$ に次数 $D-d+1$ の全単項式を掛け、
次数 $D$ の単項式基底で並べたMacaulay行列を作る。
指定した行と列による正方小行列について、次の値を得る。
\[
\begin{array}{c|c|c|c|c|c|c}
 \text{成分}&d&D&\text{行列の大きさ}&
 \text{小行列の次数}&v_{101}(\det)&\text{保証精度}\\\hline
 Z_1&6&13&135\times105&102&48&256\\
 Z_3&8&19&273\times210&205&63&256
\end{array}
\]
行列の全成分と小行列式はHensel座標の整係数多項式であり、分母を導入しない。
従って入力精度 $101^{256}$ が小行列式にもそのまま保証される。
表の付値は保証精度未満なので、真の標数零のHensel点で小行列式は非零である。

$J_G=(G_x,G_y,G_z)$ とすると、表から
\[
 \dim(S/J_G)_{13}\leq3\quad(Z_1),\qquad
 \dim(S/J_G)_{19}\leq5\quad(Z_3)
\]
を得る。標数零なのでEuler公式により $G\in J_G$ であり、
$\operatorname{Proj}(S/J_G)$ は $\Gamma_B$ のJacobian特異スキームである。
補題\ref{lem:jacobian-length-bound-v26}により、その長さはそれぞれ三、五以下である。
一方、命題\ref{prop:fixed-gram-node-v26}により、既知の相異なる三、五個の固定点は
通常二重点であり、その各点におけるJacobian特異スキームの長さは一である。
従って特異スキームはこれらの被約固定点だけに一致する。

特異点が有限なので $\Gamma_B$ は幾何学的に被約である。
もし次数 $a$ と $d-a$ の二つの正次数曲線に分けられれば、
B\'ezoutの定理により、その交差数は
\[
 a(d-a)\geq d-1
\]
である。全特異点が通常二重点なので、異なる成分の交点は横断的で、
各点の交差数は一である。
従って少なくとも $d-1$ 個の特異点が必要になるが、
ここでは $3<5$、$5<7$ なので矛盾する。
よって $\Gamma_B$ は幾何学的に既約である。
正規化の種数は平面曲線の算術種数から通常二重点数を引き、
$10-3=7$、$21-5=16$ となる。

以上の非消滅条件は、既知の固定点が相異なる通常二重点である開部分で開条件を定める。
同じHensel点が小さいGram曲線の滑らかさも満たすので、
元の固定成分内に求める非空開部分を得る。
最後に、固定点は $C$ の外にあり、これ以外に $\Gamma_B$ の特異点はない。
命題\ref{prop:complementary-gram-contact-v25}の
$\Gamma_B|_C=2D_{\rm mov}$ により、
非固定点での局所交差数は二となる。
\end{proof}

\paragraph{再現入口.}
\nolinkurl{complementary_gram_nodes_v26.py} と
\nolinkurl{complementary_gram_nodes_v26_certificate.json} が、
小行列の行番号・列番号、行列式の付値と正規化した剰余、Gram多項式のhash、
既存Hensel証明書との同一性を記録する。
既定の実行は読取り専用照合である。
この計算は標数零の特異スキームの上界を証明するものであり、
有限体の算術Galois群から幾何学的monodromyを推定するものではない。

\begin{remark}[補助曲線の対称行列表示も線束を定める]
\label{rem:complementary-nodal-prym-v26}
上の開部分では $Q_B=BAB^t$ の階数は $\Gamma_B$ の全点で $q-1$ である。
固定点では命題\ref{prop:fixed-gram-node-v26}による。
それ以外で階数が $q-2$ 以下なら随伴行列が零となり、余分な特異点が生じてしまう。
従って
\[
 0\longrightarrow\mathcal O_{\mathbb P^2}(-2)^q
 \xrightarrow{Q_B}\mathcal O_{\mathbb P^2}^q
 \longrightarrow j_*\mathcal L_B\longrightarrow0
\]
の余核は、節点を含めて $\Gamma_B$ 上の線束である。
Cartier因子に対する双対性と対称性を用いる同じ証明により、
\[
 \eta_B:=\mathcal L_B(1-q),\qquad
 \eta_B^{\otimes2}\simeq\mathcal O_{\Gamma_B},\qquad
 H^0(\Gamma_B,\eta_B)=0
\]
を得る。従って補助節点曲線にも非自明な二捩れ線束が付随し、
連結な有限エタール二重被覆を定める。
ここでは、この線束の正規化への引き戻しが非自明であることや、
小さいGram曲線の二重被覆との同一視までは主張しない。
\end{remark}

\subsection{三点固定部分の幾何学的monodromy}
\label{subsec:z1-fixed-monodromy-v26}

\begin{proposition}[固定三点への射影は全対称群]
\label{prop:z1-fixed-monodromy-v26}
$Z_1$ の被約二十一点配置上の幾何学的monodromyを
\[
 G_{Z_1}\subset
 \mathfrak S_3\times\bigl((\mathbb Z/2\mathbb Z)^9\rtimes\mathfrak S_9\bigr)
\]
と書くと、第一因子への射影は $\mathfrak S_3$ 全体である。
この命題は $\mathfrak S_3\times\{1\}\subset G_{Z_1}$ を主張しない。
\end{proposition}

\begin{proof}
ここで monodromy の基底は指定反対合を持つ枠付き成分であり、粗有理商ではない。
固定点を定める $B$ の退化スキームは、基点 $B_-$ で既に有限被約である。
正固有行を $u_{12},u_{13},u_{23}$、負固有行を $u_{11},u_{22},u_{33}$ とする順序では、
基点の負ブロックは
\[
 B_-(x,y,z)=
 \begin{pmatrix}
 0&0&0&y&0&z&0\\
 0&x&0&0&z&0&0\\
 x&0&y&0&0&0&0
 \end{pmatrix}.
\]
三行の台となる列集合は互いに交わらないので、退化点は三つの座標点だけである。
例えば $x\ne0$ の座標点では、二つの三次小行列式が $x^2y$ と $x^2z$ なので、
局所イデアルは $(y,z)$ となる。他の二点も同様である。
実際、三次小行列式のイデアルは
$(x,y,z)(xy,xz,yz)$ で、その飽和は $(xy,xz,yz)$ である。

従って $Z_1$ の滑らかな開部分 $W$ を、$B_-$ を含み、
$B$ の退化スキーム $\mathcal T\to W$ が有限エタールな三重被覆となるように取れる。
この延長の平坦性も確認できる。固有性から基点近傍で有限となり、
一般階数は三、被約な特殊ファイバーの次元も三である。
正則局所整域上で Nakayama の補題による $R^3$ から有限代数への全射を取ると、
核は自由加群内の階数零の部分加群なので零である。
従って有限自由階数三であり、被約な特殊ファイバーから近傍でエタールとなる。
以下の枠変更群の基点軌道も $W$ に含めてよい。
その軌道では退化スキームは上の三点の射影変換であり、
滑らかさと有限エタール性が保たれるためである。

$J_0=\operatorname{diag}(-1,1,1,-1,1,-1)$ を元の中央矢順序で書く。
同じ平面座標変更を左右に施し、中央矢の基底変更を $J_0$ の中心化群に取ると、
指定反対合を保つ連結群
\[
 \mathrm{GL}_3\times
 \bigl(\mathrm{GL}_3\times\mathrm{GL}_3\bigr)
\]
が作用する。関係枠を保つ表示では、さらに連結群 $\mathrm{GL}_7$ を付けてもよい。
この群は $Z_1$ を保つ。

三巡回置換 $\sigma=(123)$ に対する
$(\sigma,\operatorname{Sym}^2\sigma)$ と、互換 $\rho=(12)$ に対する
$(\rho,J_0\operatorname{Sym}^2\rho)$ は、ともに $B_-$ を固定する。
これを左関係の列 $r_0,\ldots,r_6$ で確認すると、像はそれぞれ
\[
 (r_3,r_2,r_5,r_4,r_0,r_1,r_6),\qquad
 (-r_2,-r_3,-r_0,-r_1,-r_5,-r_4,r_6)
\]
となる。中央作用はいずれも $J_0$ と可換であり、右関係も同時に保たれる。
関係pairingについても、対応する七次行列 $M$ は $MAM^t=A$ を満たす。

複素数上でこの連結枠変更群の恒等元から各安定化元まで経路を取り、
$B_-$ に作用させると、$W$ 内の基点軌道上の閉路を得る。
その三点被覆への作用は座標点の三巡回置換と互換である。
従って $\mathcal T\to W$ の幾何学的monodromyは $\mathfrak S_3$ である。

最後に、全二十一点が有限エタールになる非空開部分を $U\subset W$ とする。
滑らかな既約基底の稠密開部分への制限で、有限エタール被覆のmonodromyは縮まらない。
例えば被覆の連結Galois閉包は既約であり、その非空開部分も連結なことから従う。
よって $\mathcal T|_U$ のmonodromyも $\mathfrak S_3$ である。
これは全二十一点のmonodromyの固定三点への射影に一致する。
\end{proof}

\paragraph{再現入口.}
\nolinkurl{z1_fixed_monodromy_v26.py} と
\nolinkurl{z1_fixed_monodromy_v26_certificate.json} が、
三次小行列式の単項式生成元、二つの中央行列と関係行列、
$J_0$ との可換性、関係およびpairingの保存を有理数上で厳密に確認する。
幾何学的monodromyへの移行は、上の連結群軌道と稠密開制限の証明による。


\section{検証と残る問題}
付属の厳密な検証は、同じ既存成分上での滑らかさ、階数七のGram復元、
九余因子の標準形、標準十八点の分解、七次被覆の単純分岐、
収縮二次曲線の一次延長障害、残余六次曲線の滑らかさ、
三十節点の平滑化、および動く基点を追った六点の二次jetを対象とする。
幾何学的モノドロミー $S_7$ は、源の既約性による推移性、素数次数による原始性、
および一つの単純分岐に対する局所互換を組み合わせて証明する。
有限体での多項式の分解型だけから幾何学的群を推定するものではない。

九余因子方程式では、整数基点を通る既約成分を指定している。
全零点集合の既約性、大域的有理性、分岐曲線の特異点型、
分岐曲線からの元のPrym対の復元、およびPrym像全体の双有理型は未決である。
通常尖点七十二個・節点百四十四個という個数は、特異点がその二種類だけと証明された場合の帰結であり、
本稿は正規化写像の重み付き微分零因子と $\delta$ の総和を確定する。

\clearpage
\begingroup\small
\begin{thebibliography}{99}
\setlength{\itemsep}{2pt}
\setlength{\parskip}{0pt}
\bibitem{PaperII}
\emph{Fano72型関係代数の二十一点Torelliと反対代数}, 本研究の先行稿、第26版、2026年。
\nolinkurl{fano72_point_configurations_ja_v26.pdf}.

\bibitem{BeauvilleDeterminantal}
A.~Beauville, \emph{Determinantal hypersurfaces}, Michigan Math. J.\ \textbf{48} (2000), 39--64.
\url{https://arxiv.org/abs/math/9910030}.

\bibitem{BoehningBothmerPrym}
C.~B\"ohning and H.-C.~Graf von Bothmer,
\emph{On stable rationality of some conic bundles and moduli spaces of Prym curves},
Comment. Math. Helv.\ \textbf{93} (2018), 133--155.
\url{https://doi.org/10.4171/CMH/431}.

\bibitem{HartshorneNoether}
R.~Hartshorne, \emph{Generalized divisors on Gorenstein curves and a theorem of Noether},
J. Math. Kyoto Univ.\ \textbf{26} (1986), no.~3, 375--386.
\url{https://doi.org/10.1215/kjm/1250520873}.

\bibitem{CoppensKatoPlaneSystems}
M.~Coppens and T.~Kato, \emph{Non-trivial linear systems on smooth plane curves},
Math. Nachr.\ \textbf{166} (1994), 71--82.
\url{https://doi.org/10.1002/mana.19941660106}.
\url{https://arxiv.org/abs/alg-geom/9301003}.

\bibitem{EisenbudSidman}
D.~Eisenbud, with a chapter by J.~Sidman,
\emph{Lectures on the geometry of syzygies},
in Trends in Commutative Algebra, MSRI Publications \textbf{51} (2004), 115--152.
\url{https://library.slmath.org/books/Book51/files/04eisenbud.pdf}.

\bibitem{PaperV}
\emph{四次Prymデータに付随する線織面と種数二十K3の境界周期}, 本研究の公表稿V、第30版、2026年。
\nolinkurl{self_opposite_prym_fibrations_ja_v30.pdf}.
\end{thebibliography}
\endgroup
\end{document}
